% main.tex - 中文线性代数讲义
% 主文档 - 使用 XeLaTeX 编译
%
% 线性代数: 线性代数的本质 (Essence of Linear Algebra)
% 基于 3Blue1Brown 视频系列
%
% 编译方式: xelatex main.tex (运行 2 次以生成目录和引用)

% preamble.tex - 中文线性代数讲义
% XeLaTeX 配置 (ctex 中文支持)
% 编译方式: xelatex main.tex (运行 2 次以生成目录和引用)

\documentclass[12pt,a4paper,openany,oneside]{book}

% === 中文支持 (必须最先加载, 需要 XeLaTeX) ===
\usepackage[UTF8, scheme=plain]{ctex}

% === 页面设置 ===
\usepackage{geometry}
\geometry{
    top=2.5cm,
    bottom=2.5cm,
    left=2.5cm,
    right=2.5cm,
    headheight=14pt
}

% === 数学宏包 ===
\usepackage{amsmath}
\usepackage{amssymb}
\usepackage{amsthm}
\usepackage{mathtools}
\usepackage{bm}

% === 图形与绘图 ===
\usepackage{graphicx}
\usepackage{tikz}
\usepackage{pgfplots}
\pgfplotsset{compat=1.18}
\usetikzlibrary{arrows.meta, calc, positioning, decorations.pathreplacing, patterns}

% === 颜色与盒子 ===
\usepackage{xcolor}
\usepackage{tcolorbox}
\tcbuselibrary{skins,breakable,theorems}

% === 表格与列表 ===
\usepackage{booktabs}
\usepackage{array}
\usepackage{enumitem}

% === 超链接 (建议放在较后位置) ===
\usepackage[colorlinks=true,linkcolor=blue!70!black,urlcolor=blue!70!black,citecolor=green!50!black]{hyperref}

% === 定理环境 (中文) ===
\theoremstyle{definition}
\newtheorem{definition}{定义}[chapter]
\newtheorem{example}{例}[chapter]
\newtheorem{exercise}{练习}[chapter]
\newtheorem{problem}{习题}[chapter]

\theoremstyle{plain}
\newtheorem{theorem}{定理}[chapter]
\newtheorem{lemma}{引理}[chapter]
\newtheorem{proposition}{命题}[chapter]
\newtheorem{corollary}{推论}[chapter]

\theoremstyle{remark}
\newtheorem{remark}{注}[chapter]
\newtheorem{note}{备注}[chapter]

% === 自定义 tcolorbox 环境 (中文标题) ===

% 几何直观盒子
\newtcolorbox{intuition}{
    enhanced,
    colback=blue!5,
    colframe=blue!75!black,
    fonttitle=\bfseries,
    title={几何直观},
    breakable,
    left=5mm,
    right=5mm,
    top=3mm,
    bottom=3mm,
    arc=2mm,
    boxrule=1pt
}

% 实际应用盒子
\newtcolorbox{practical}{
    enhanced,
    colback=green!5,
    colframe=green!60!black,
    fonttitle=\bfseries,
    title={实际应用},
    breakable,
    left=5mm,
    right=5mm,
    top=3mm,
    bottom=3mm,
    arc=2mm,
    boxrule=1pt
}

% 注意盒子
\newtcolorbox{warning}{
    enhanced,
    colback=red!5,
    colframe=red!75!black,
    fonttitle=\bfseries,
    title={注意},
    breakable,
    left=5mm,
    right=5mm,
    top=3mm,
    bottom=3mm,
    arc=2mm,
    boxrule=1pt
}

% 重点盒子
\newtcolorbox{important}{
    enhanced,
    colback=orange!5,
    colframe=orange!75!black,
    fonttitle=\bfseries,
    title={重点},
    breakable,
    left=5mm,
    right=5mm,
    top=3mm,
    bottom=3mm,
    arc=2mm,
    boxrule=1pt
}

% 本章小结盒子
\newtcolorbox{summary}{
    enhanced,
    colback=purple!5,
    colframe=purple!75!black,
    fonttitle=\bfseries,
    title={本章小结},
    breakable,
    left=5mm,
    right=5mm,
    top=3mm,
    bottom=3mm,
    arc=2mm,
    boxrule=1pt
}

% === 页眉页脚 ===
\usepackage{fancyhdr}
\pagestyle{fancy}
\fancyhf{}
\fancyhead[L]{\leftmark}
\fancyhead[R]{\rightmark}
\fancyfoot[C]{\thepage}
\renewcommand{\headrulewidth}{0.4pt}
\renewcommand{\footrulewidth}{0pt}

% === 行距 ===
\usepackage{setspace}
\onehalfspacing

% === 摘要环境 (book 文档类默认没有 abstract 环境, 自定义为缩进小字块) ===
\newenvironment{abstract}{%
  \small
  \begin{quotation}\noindent\ignorespaces
}{%
  \end{quotation}
}

% === 自定义数学命令 ===
% commands.tex - 线性代数自定义数学命令
% 中文线性代数讲义 (与英文版一致, 另补英文版漏定义的 \vp 等)

% === Vector Notation ===
% Bold vectors (column vectors)
\newcommand{\vect}[1]{\mathbf{#1}}
\newcommand{\vv}{\vect{v}}
\newcommand{\vw}{\vect{w}}
\newcommand{\vu}{\vect{u}}
\newcommand{\vx}{\vect{x}}
\newcommand{\vy}{\vect{y}}
\newcommand{\vz}{\vect{z}}
\newcommand{\va}{\vect{a}}
\newcommand{\vb}{\vect{b}}
\newcommand{\vc}{\vect{c}}
\newcommand{\vd}{\vect{d}}
\newcommand{\ve}{\vect{e}}
\newcommand{\vf}{\vect{f}}
\newcommand{\vg}{\vect{g}}
\newcommand{\vh}{\vect{h}}
\newcommand{\vl}{\vect{l}}
\newcommand{\vm}{\vect{m}}
\newcommand{\vn}{\vect{n}}
\newcommand{\vo}{\vect{o}}
\newcommand{\vp}{\vect{p}}
\newcommand{\vq}{\vect{q}}
\newcommand{\vr}{\vect{r}}
\newcommand{\vs}{\vect{s}}
\newcommand{\vt}{\vect{t}}

% Basis vectors (with hat notation)
\newcommand{\vi}{\hat{\imath}}
\newcommand{\vj}{\hat{\jmath}}
\newcommand{\vk}{\hat{k}}

% Zero vector
\newcommand{\vzero}{\vect{0}}

% === Matrix Notation ===
\newcommand{\mat}[1]{\mathbf{#1}}
\newcommand{\mA}{\mat{A}}
\newcommand{\mB}{\mat{B}}
\newcommand{\mC}{\mat{C}}
\newcommand{\mD}{\mat{D}}
\newcommand{\mE}{\mat{E}}
\newcommand{\mF}{\mat{F}}
\newcommand{\mG}{\mat{G}}
\newcommand{\mH}{\mat{H}}
\newcommand{\mI}{\mat{I}}
\newcommand{\mJ}{\mat{J}}
\newcommand{\mK}{\mat{K}}
\newcommand{\mL}{\mat{L}}
\newcommand{\mM}{\mat{M}}
\newcommand{\mN}{\mat{N}}
\newcommand{\mO}{\mat{O}}
\newcommand{\mP}{\mat{P}}
\newcommand{\mQ}{\mat{Q}}
\newcommand{\mR}{\mat{R}}
\newcommand{\mS}{\mat{S}}
\newcommand{\mT}{\mat{T}}
\newcommand{\mU}{\mat{U}}
\newcommand{\mV}{\mat{V}}
\newcommand{\mW}{\mat{W}}
\newcommand{\mX}{\mat{X}}
\newcommand{\mY}{\mat{Y}}
\newcommand{\mZ}{\mat{Z}}

% Matrix with brackets
\newcommand{\bmat}[1]{\begin{bmatrix}#1\end{bmatrix}}
\newcommand{\pmat}[1]{\begin{pmatrix}#1\end{pmatrix}}
\newcommand{\vmat}[1]{\begin{vmatrix}#1\end{vmatrix}}

% === Common Operations ===
% Transpose
\newcommand{\trans}{^{\mathsf{T}}}

% Inverse
\newcommand{\inv}{^{-1}}

% Norm and absolute value
\newcommand{\norm}[1]{\left\| #1 \right\|}
\newcommand{\abs}[1]{\left| #1 \right|}

% === Important Spaces ===
\newcommand{\R}{\mathbb{R}}
\newcommand{\C}{\mathbb{C}}
\newcommand{\N}{\mathbb{N}}
\newcommand{\Z}{\mathbb{Z}}
\newcommand{\Q}{\mathbb{Q}}

% Rn space
\newcommand{\Rtwo}{\mathbb{R}^2}
\newcommand{\Rthree}{\mathbb{R}^3}

% === Linear Algebra Operators ===
\DeclareMathOperator{\spn}{span}
\DeclareMathOperator{\Span}{Span}
\DeclareMathOperator{\Null}{Null}
\DeclareMathOperator{\Col}{Col}
\DeclareMathOperator{\Row}{Row}
\DeclareMathOperator{\rank}{rank}
\DeclareMathOperator{\tr}{tr}
\DeclareMathOperator{\proj}{proj}

% === Common Matrices ===
% 2x2 matrix shorthand
\newcommand{\twomat}[4]{\begin{bmatrix} #1 & #2 \\ #3 & #4 \end{bmatrix}}

% 3x3 matrix shorthand
\newcommand{\threemat}[9]{\begin{bmatrix} #1 & #2 & #3 \\ #4 & #5 & #6 \\ #7 & #8 & #9 \end{bmatrix}}

% 2D column vector
\newcommand{\twovec}[2]{\begin{bmatrix} #1 \\ #2 \end{bmatrix}}

% 3D column vector
\newcommand{\threevec}[3]{\begin{bmatrix} #1 \\ #2 \\ #3 \end{bmatrix}}

% === Change of Basis ===
\newcommand{\coords}[2]{[#1]_{#2}}

% === System of Equations ===
\newcommand{\system}[1]{\left\{\begin{aligned} #1 \end{aligned}\right.}

% === Emphasizing ===
\newcommand{\vocab}[1]{\textbf{\textcolor{blue!70!black}{#1}}}


% === TikZ 线性代数绘图样式 ===
\tikzset{
    vector/.style={-{Stealth[length=3mm]}, thick},
    axis/.style={-{Stealth[length=2mm]}, gray},
    grid/.style={very thin, gray!30},
    transformed/.style={blue!70},
    original/.style={black},
    basis/.style={red!70!black, very thick},
    point/.style={circle, fill, inner sep=1.5pt}
}


\begin{document}

% === 封面 ===
\begin{titlepage}
\centering
\vspace*{1cm}

{\Huge\bfseries 线性代数}\\[0.5cm]
{\LARGE\bfseries 线性代数的本质}\\[0.3cm]
{\large Essence of Linear Algebra}\\[2cm]

{\Large 综合讲义}\\[0.5cm]
{\large 强调几何直观与实际应用}\\[2cm]

\begin{tikzpicture}[scale=0.6]
    % 背景网格
    \draw[grid] (-3,-3) grid (3,3);

    % 坐标轴
    \draw[axis] (-3.5,0) -- (3.5,0);
    \draw[axis] (0,-3.5) -- (0,3.5);

    % 基向量
    \draw[basis, -{Stealth[length=4mm]}] (0,0) -- (1,0) node[below right] {$\vi$};
    \draw[basis, -{Stealth[length=4mm]}] (0,0) -- (0,1) node[above left] {$\vj$};

    % 一个变换后的向量
    \draw[vector, blue!70, very thick] (0,0) -- (2,1.5) node[above right] {$\vv$};

    % 另一个向量
    \draw[vector, green!60!black, very thick] (0,0) -- (-1,2) node[above left] {$\vw$};
\end{tikzpicture}\\[2cm]

{\large 基于 Grant Sanderson 的视频系列}\\[0.3cm]
{\Large\textbf{线性代数的本质}}\\[0.3cm]
{\large Essence of Linear Algebra}\\[0.3cm]
{\small YouTube 频道: 3Blue1Brown}\\[1.5cm]

{\small \today}

\end{titlepage}

% === 致谢页 ===
\begin{titlepage}
\centering
\vspace*{3cm}

{\Large\bfseries 致谢}

\vspace{2cm}

{\large 特别感谢原英文/波斯语版贡献者:}

\vspace{1cm}

Ebrahim Khosravani\\[0.25cm]
Mohammadhasan Shiri\\[0.25cm]
Zahra Amiri\\[0.25cm]
Amin Azizi\\[0.25cm]
Bita Qorbani\\[0.25cm]
Yasamin Yazdani\\[0.25cm]
Sepehr Talaei

\vspace{2cm}

{\small 基于 3Blue1Brown 的``线性代数的本质 (Essence of Linear Algebra)''视频系列}

\vspace{1cm}

{\small 谢里夫理工大学 (Sharif University of Technology) 物理学科学协会}

\vspace{1cm}
{\small 中文版由英文版翻译而来, 仅供学习交流使用}
\end{titlepage}

% === 目录 ===
\tableofcontents

% === 前言 ===
\chapter*{前言}
\addcontentsline{toc}{chapter}{前言}

本讲义基于 Grant Sanderson 在 YouTube 频道 3Blue1Brown 上发布的著名视频系列``线性代数的本质 (Essence of Linear Algebra)''。该系列的主要目标是帮助读者建立对线性代数概念的直观、几何化理解。

\section*{本讲义的特点}
\begin{itemize}
    \item \textbf{几何直观 (Geometric Intuition):} 强调对概念的可视化理解, 而不仅仅是计算
    \item \textbf{严格定义 (Rigorous Definitions):} 标准的数学表述
    \item \textbf{实际应用 (Practical Applications):} 在物理、计算机科学与工程中的应用
    \item \textbf{练习 (Exercises):} 类型多样的习题以巩固学习
\end{itemize}

\section*{预备知识}
\begin{itemize}
    \item 基本数学知识 (高中代数)
    \item 笛卡尔坐标系 (Cartesian coordinate system) 的理解
    \item 对深入理解数学的渴望!
\end{itemize}

\section*{如何使用本讲义}
每章包含以下部分:
\begin{description}
    \item[定义:] 关键概念与数学符号
    \item[几何直观:] 可视化解释 (蓝色框)
    \item[实际应用:] 现实世界中的例子 (绿色框)
    \item[练习:] 供实践的习题
\end{description}

\begin{flushright}
祝线性代数学习愉快!
\end{flushright}

% === 各章 ===

% 第 1 章: 向量简介
% lecture01.tex - 向量简介
% 第 1 章: 向量简介 (Introduction to Vectors)

\chapter{向量简介}
\label{ch:vectors}

\begin{abstract}
向量 (vector) 是线性代数的基本构件。在本章中, 我们将从三个不同的视角来认识向量: 物理视角、计算机科学视角和数学视角。我们还将学习向量加法 (vector addition) 与标量乘法 (scalar multiplication) 这两种基本运算。
\end{abstract}

% ============================================
\section{看待向量的三种视角}
% ============================================

向量 (vector) 这个概念在不同学科中有着不同的含义。理解这三种视角有助于我们形成对向量的完整认识。

% -------------------------------------------
\subsection{物理视角}
% -------------------------------------------

\begin{definition}[物理视角下的向量]
向量是空间中的一支\vocab{箭头 (arrow)}, 由两个属性决定:
\begin{enumerate}
    \item \textbf{长度 (length/magnitude)}
    \item \textbf{方向 (direction)}
\end{enumerate}
只要这两个属性不变, 向量可以在空间中任意平移, 仍被视为同一个向量。
\end{definition}

\begin{intuition}
在纸上画一支箭头, 现在把它平移到另一个位置, 不旋转也不改变大小。从物理学家的角度看, 这仍然是同一个向量!

例如, 苹果所受的重力始终竖直向下, 大小恒定---无论苹果在房间的哪个位置。
\end{intuition}

\begin{center}
\begin{tikzpicture}[scale=1.2]
    % 画出同一个向量的多个副本
    \draw[vector, blue!70, very thick] (0,0) -- (2,1) node[midway, above] {$\vv$};
    \draw[vector, blue!70, very thick] (3,1) -- (5,2) node[midway, above] {$\vv$};
    \draw[vector, blue!70, very thick] (-1,2) -- (1,3) node[midway, above] {$\vv$};

    % 注释
    \node at (2,-0.5) {\small 这些都是同一个向量};
\end{tikzpicture}
\end{center}

% -------------------------------------------
\subsection{计算机科学视角}
% -------------------------------------------

\begin{definition}[计算机科学视角下的向量]
向量是\vocab{有序数表 (ordered list of numbers)}。例如:
\[
\vv = \twovec{3}{-2}
\]
在这种视角下, ``向量 (vector)''几乎就是``数表 (list)''的同义词。
\end{definition}

\begin{practical}
\textbf{例: 房价模型}

假设我们想用两个特征来描述房子:
\begin{itemize}
    \item 面积 (平方米)
    \item 价格 (千美元)
\end{itemize}

每栋房子就是一个二维向量:
\[
\text{House}_1 = \twovec{85}{250}, \quad
\text{House}_2 = \twovec{120}{420}, \quad
\text{House}_3 = \twovec{65}{180}
\]

\textbf{重点:} 数字的顺序很重要! $\twovec{85}{250}$ 与 $\twovec{250}{85}$ 是不同的向量。
\end{practical}

% -------------------------------------------
\subsection{数学视角}
% -------------------------------------------

\begin{definition}[数学视角下的向量]
数学家对向量作了抽象定义: 向量是任何可以进行以下两种运算的对象:
\begin{enumerate}
    \item \textbf{两个向量相加 (add two vectors)}
    \item \textbf{向量乘以一个数, 即标量乘法 (scalar multiplication)}
\end{enumerate}
我们将在最后一章 (抽象向量空间) 中详细探讨这一定义。
\end{definition}

\begin{important}
在本课程中, 我们把向量可视化为\vocab{坐标系中以原点 (origin) 为起点}的一支\textbf{箭头}。这种几何视角有助于我们更深刻地理解各种概念。
\end{important}

% ============================================
\section{坐标系与向量的表示}
% ============================================

\begin{definition}[二维笛卡尔坐标系 (Cartesian coordinate system)]
一个坐标系由以下部分组成:
\begin{itemize}
    \item 水平轴: \vocab{$x$ 轴 ($x$-axis)}
    \item 竖直轴: \vocab{$y$ 轴 ($y$-axis)}
    \item 两轴交点: \vocab{原点 (origin)}
\end{itemize}
\end{definition}

\begin{center}
\begin{tikzpicture}[scale=1.5]
    % 网格
    \draw[grid] (-0.5,-0.5) grid (4.5,3.5);

    % 坐标轴
    \draw[axis, thick] (-0.5,0) -- (4.5,0) node[right] {$x$};
    \draw[axis, thick] (0,-0.5) -- (0,3.5) node[above] {$y$};

    % 原点标记
    \node[below left] at (0,0) {O};

    % 刻度
    \foreach \x in {1,2,3,4} {
        \draw (\x,0.1) -- (\x,-0.1) node[below] {\small $\x$};
    }
    \foreach \y in {1,2,3} {
        \draw (0.1,\y) -- (-0.1,\y) node[left] {\small $\y$};
    }

    % 向量
    \draw[vector, blue!70, ultra thick] (0,0) -- (3,2);
    \node[above right, blue!70!black] at (3,2) {$\vv = \twovec{3}{2}$};

    % 显示坐标的虚线
    \draw[dashed, red!70] (3,0) -- (3,2);
    \draw[dashed, red!70] (0,2) -- (3,2);

    % 坐标标记
    \node[below, red!70!black] at (3,0) {$3$};
    \node[left, red!70!black] at (0,2) {$2$};
\end{tikzpicture}
\end{center}

\begin{definition}[向量的坐标 (coordinates)]
向量的\vocab{坐标}是一对数字, 描述了从原点走到向量箭头终点的方法:
\begin{itemize}
    \item 第一个数: 沿 $x$ 轴走多远 (向右为正, 向左为负)
    \item 第二个数: 沿 $y$ 轴走多远 (向上为正, 向下为负)
\end{itemize}
\end{definition}

\begin{intuition}
坐标 $\twovec{3}{2}$ 的意思是:
\begin{enumerate}
    \item 从原点出发, 先向右走 3 个单位
    \item 再向上走 2 个单位
    \item 向量的箭头终点就在这里!
\end{enumerate}
\end{intuition}

% ============================================
\section{三维空间}
% ============================================

\begin{definition}[三维向量 (three-dimensional vector)]
在三维空间中, 增加了一条同时垂直于 $x$ 轴和 $y$ 轴的\vocab{$z$ 轴 ($z$-axis)}。每个向量由三个数确定:
\[
\vv = \threevec{a}{b}{c}
\]
\end{definition}

\begin{center}
\begin{tikzpicture}[scale=1.2]
    % 三维坐标轴
    \draw[axis, thick] (0,0) -- (3,0) node[right] {$x$};
    \draw[axis, thick] (0,0) -- (0,3) node[above] {$z$};
    \draw[axis, thick] (0,0) -- (-1.5,-1) node[below left] {$y$};

    % 向量
    \draw[vector, blue!70, ultra thick] (0,0) -- (2,1.5);
    \draw[vector, green!60!black, ultra thick] (0,0) -- (-0.8,-0.5);

    % 标记
    \node[above right, blue!70!black] at (2,1.5) {$\threevec{2}{0}{1.5}$};
\end{tikzpicture}
\end{center}

% ============================================
\section{向量加法}
% ============================================

\begin{definition}[两向量相加---几何方法]
将两个向量 $\va$ 与 $\vb$ 相加:
\begin{enumerate}
    \item 先画出向量 $\va$
    \item 把向量 $\vb$ 平移, 使它的起点落在 $\va$ 的终点上
    \item 和向量就是从原点 ($\va$ 的起点) 指向 $\vb$ 终点的向量
\end{enumerate}
\end{definition}

\begin{center}
\begin{tikzpicture}[scale=1.3]
    % 网格
    \draw[grid] (-0.5,-0.5) grid (5.5,3.5);

    % 坐标轴
    \draw[axis] (-0.5,0) -- (5.5,0);
    \draw[axis] (0,-0.5) -- (0,3.5);

    % 向量 a
    \draw[vector, red!70, very thick] (0,0) -- (2,1) node[midway, below] {$\va$};

    % 平移后的向量 b
    \draw[vector, blue!70, very thick] (2,1) -- (4,3) node[midway, right] {$\vb$};

    % 和向量
    \draw[vector, purple!70, ultra thick] (0,0) -- (4,3) node[midway, above left] {$\va + \vb$};

    % 原始的 b (虚线)
    \draw[vector, blue!40, dashed] (0,0) -- (2,2) node[midway, left] {\small $\vb$};
\end{tikzpicture}
\end{center}

\begin{definition}[两向量相加---代数方法]
若 $\va = \twovec{a_1}{a_2}$, $\vb = \twovec{b_1}{b_2}$:
\[
\va + \vb = \twovec{a_1 + b_1}{a_2 + b_2}
\]
即把对应分量分别相加。
\end{definition}

\begin{example}
\[
\twovec{1}{2} + \twovec{3}{-1} = \twovec{1+3}{2+(-1)} = \twovec{4}{1}
\]
\end{example}

\begin{intuition}
向量加法可以理解为\textbf{沿路径行走}:
\begin{itemize}
    \item 先沿向量 $\va$ 走一段
    \item 再沿向量 $\vb$ 走一段
    \item 最终效果与直接沿 $\va + \vb$ 走相同
\end{itemize}

就像走路: 先向右走 2 步, 再向右走 5 步, 等同于直接向右走 7 步。
\end{intuition}

\begin{practical}
\textbf{例: 风中的飞机速度}

一架飞机以 $\twovec{500}{0}$ km/h 向东飞行 ($x$ 正方向)。风以 $\twovec{0}{50}$ km/h 向北吹 ($y$ 正方向)。

飞机相对地面的实际速度为:
\[
\vv_{\text{actual}} = \twovec{500}{0} + \twovec{0}{50} = \twovec{500}{50}
\]

飞机既向东飞, 又略微向北偏移。
\end{practical}

% ============================================
\section{标量乘法}
% ============================================

\begin{definition}[数乘向量]
用一个数 (标量, scalar) $c$ 乘以向量 $\vv$, 就是把向量\vocab{缩放 (scale)} $c$ 倍:
\[
c \cdot \twovec{v_1}{v_2} = \twovec{c \cdot v_1}{c \cdot v_2}
\]
\end{definition}

\begin{center}
\begin{tikzpicture}[scale=1]
    % 原向量
    \draw[vector, blue!70, very thick] (0,0) -- (1.5,1) node[right] {$\vv$};

    % 缩放后的向量
    \draw[vector, green!60!black, very thick] (3,0) -- (6,2) node[right] {$2\vv$};
    \draw[vector, orange!70, very thick] (8,0) -- (8.75,0.5) node[right] {$\frac{1}{2}\vv$};
    \draw[vector, red!70, very thick] (11,0) -- (8.3,-1.8) node[below] {$-1.8\vv$};

    % 下方标记
    \node at (0.75,-0.5) {\small 原向量};
    \node at (4.5,-0.5) {\small 拉伸};
    \node at (8.4,-0.5) {\small 压缩};
    \node at (9.7,-2.3) {\small 反向并拉伸};
\end{tikzpicture}
\end{center}

\begin{important}
标量乘法的效果:
\begin{itemize}
    \item $c > 1$: 向量被\textbf{拉伸 (stretched)}
    \item $0 < c < 1$: 向量被\textbf{压缩 (compressed)}
    \item $c < 0$: 向量先\textbf{反向 (reversed)}, 再缩放
    \item $c = 0$: 向量变为零向量
    \item $c = 1$: 向量不变
\end{itemize}
\end{important}

\begin{warning}
若 $c < 0$, 向量不仅大小改变, \textbf{方向也会反转}!
\end{warning}

\begin{example}
若 $\vv = \twovec{4}{-2}$:
\begin{align*}
2\vv &= \twovec{8}{-4} & &\text{(变为两倍)} \\
-\vv &= \twovec{-4}{2} & &\text{(反向)} \\
\frac{1}{2}\vv &= \twovec{2}{-1} & &\text{(变为一半)}
\end{align*}
\end{example}

\begin{definition}[标量 (scalar)]
在线性代数中, 与向量相乘的数称为\vocab{标量}。这个名字来自动词``缩放 (to scale)''。``标量 (scalar)''本质上就是``数''的同义词。
\end{definition}

% ============================================
\section{基向量}
% ============================================

\begin{definition}[$\Rtwo$ 中的标准基向量 (standard basis vectors)]
两个标准基向量为:
\[
\vi = \twovec{1}{0} \quad \text{和} \quad \vj = \twovec{0}{1}
\]
\end{definition}

\begin{center}
\begin{tikzpicture}[scale=2]
    % 网格
    \draw[grid] (-0.5,-0.5) grid (2.5,2.5);

    % 坐标轴
    \draw[axis] (-0.5,0) -- (2.5,0) node[right] {$x$};
    \draw[axis] (0,-0.5) -- (0,2.5) node[above] {$y$};

    % 基向量
    \draw[basis, -{Stealth[length=4mm]}] (0,0) -- (1,0) node[below right] {$\vi$};
    \draw[basis, -{Stealth[length=4mm]}] (0,0) -- (0,1) node[above left] {$\vj$};

    % 一般向量
    \draw[vector, blue!70, very thick] (0,0) -- (2,1.5);
    \node[above right, blue!70!black] at (2,1.5) {$\vv = 2\vi + 1.5\vj$};

    % 分量
    \draw[dashed, orange] (0,0) -- (2,0) node[midway, below] {$2\vi$};
    \draw[dashed, orange] (2,0) -- (2,1.5) node[midway, right] {$1.5\vj$};
\end{tikzpicture}
\end{center}

\begin{theorem}
$\Rtwo$ 中的任何向量都可以写成基向量的线性组合 (linear combination):
\[
\twovec{a}{b} = a\vi + b\vj = a\twovec{1}{0} + b\twovec{0}{1}
\]
\end{theorem}

\begin{intuition}
当我们写 $\twovec{3}{2}$ 时, 实际是在说:
\begin{center}
``取 3 个 $\vi$, 再取 2 个 $\vj$, 然后把它们加起来''
\end{center}
即:
\[
\twovec{3}{2} = 3\vi + 2\vj = 3\twovec{1}{0} + 2\twovec{0}{1}
\]
\end{intuition}

% ============================================
\section{三种视角的联系}
% ============================================

\begin{summary}
\textbf{线性代数的威力}在于不同视角之间的转换:
\begin{itemize}
    \item \textbf{数据分析师:} 可以把很长的数表看作空间中的向量
    \item \textbf{物理学家:} 可以用数字来描述运动和力
    \item \textbf{图形程序员:} 可以用矩阵实现几何变换
\end{itemize}
\end{summary}

% ============================================
\section{练习}
% ============================================

\begin{exercise}
在坐标系中画出以下向量:
\[
\va = \twovec{3}{1}, \quad \vb = \twovec{-2}{4}, \quad \vc = \twovec{0}{-3}
\]
\end{exercise}

\begin{exercise}
计算以下向量的和与差:
\[
\va = \twovec{2}{5}, \quad \vb = \twovec{-1}{3}
\]
\begin{enumerate}[label=(\alph*)]
    \item $\va + \vb$
    \item $\va - \vb$
    \item $2\va + 3\vb$
\end{enumerate}
\end{exercise}

\begin{exercise}
若 $\vv = \twovec{4}{-2}$, 计算并画出以下向量:
\begin{enumerate}[label=(\alph*)]
    \item $2\vv$
    \item $-\vv$
    \item $\frac{1}{2}\vv$
    \item $-2.5\vv$
\end{enumerate}
\end{exercise}

\begin{exercise}[应用]
一艘船以 30 km/h 向北航行, 水流以 10 km/h 向东流动。船相对岸的实际速度是多少?
\end{exercise}

\begin{exercise}
证明对任何向量 $\vv$:
\[
\vv + (-\vv) = \vzero
\]
其中 $\vzero = \twovec{0}{0}$ 是零向量 (zero vector)。
\end{exercise}

\begin{exercise}
证明向量加法满足交换律 (commutative):
\[
\va + \vb = \vb + \va
\]
\end{exercise}

\begin{exercise}[挑战]
给定三点 $A(1,2)$, $B(4,6)$ 和 $C(7,2)$。证明这三点构成一个等腰三角形。

\textit{提示: 向量 $\twovec{a}{b}$ 的长度等于 $\sqrt{a^2 + b^2}$。}
\end{exercise}

\begin{problem}
在三维空间中, 给定向量 $\vv = \threevec{2}{3}{-1}$。计算:
\begin{enumerate}[label=(\alph*)]
    \item $3\vv$
    \item $\vv + \threevec{1}{-1}{2}$
    \item $-\frac{1}{2}\vv$
\end{enumerate}
\end{problem}


% 第 2 章: 线性组合、张成与基
% lecture02.tex - 线性组合、张成与基
% 第 2 章: 线性组合、张成与基

\chapter{线性组合、张成与基}
\label{ch:span-basis}

\begin{abstract}
在本章中, 我们将探讨``线性组合 (linear combination)''、``张成 (span)''与``基 (basis)''这些核心概念。这些概念是深入理解线性代数的基础, 让我们明白向量 (vector) 是如何``填满''空间的。
\end{abstract}

% ============================================
\section{坐标的新视角}
% ============================================

在上一章中, 我们把向量坐标看作``移动指南''。但还有另一种非常重要的视角。

\begin{intuition}
当我们写 $\vv = \twovec{3}{-2}$ 时, 不要只想``向右走 3 个单位, 向下走 2 个单位'', 而要这样想:
\begin{itemize}
    \item 数字 3 是一个\textbf{标量 (scalar)}, 用来缩放向量 $\vi$
    \item 数字 $-2$ 是一个\textbf{标量}, 用来缩放向量 $\vj$
    \item 最终的向量是这两个缩放后向量的\textbf{和}
\end{itemize}
\end{intuition}

\begin{center}
\begin{tikzpicture}[scale=1.3]
    % 网格
    \draw[grid] (-1.5,-2.5) grid (4.5,1.5);

    % 坐标轴
    \draw[axis] (-1.5,0) -- (4.5,0) node[right] {$x$};
    \draw[axis] (0,-2.5) -- (0,1.5) node[above] {$y$};

    % 基向量 (单位向量)
    \draw[basis, -{Stealth[length=3mm]}] (0,0) -- (1,0) node[below right] {$\vi$};
    \draw[basis, -{Stealth[length=3mm]}] (0,0) -- (0,1) node[above left] {$\vj$};

    % 缩放后的向量
    \draw[vector, orange!70, very thick] (0,0) -- (3,0) node[midway, below] {$3\vi$};
    \draw[vector, green!60!black, very thick] (3,0) -- (3,-2) node[midway, right] {$-2\vj$};

    % 结果向量
    \draw[vector, blue!70, ultra thick] (0,0) -- (3,-2) node[below right] {$\vv = 3\vi - 2\vj$};
\end{tikzpicture}
\end{center}

% ============================================
\section{线性组合}
% ============================================

\begin{definition}[线性组合 (linear combination)]
两个向量 $\vv$ 与 $\vw$ 的\vocab{线性组合}为:
\[
a\vv + b\vw
\]
其中 $a$ 与 $b$ 为任意标量。

更一般地, $n$ 个向量 $\vv_1, \vv_2, \ldots, \vv_n$ 的线性组合为:
\[
c_1\vv_1 + c_2\vv_2 + \cdots + c_n\vv_n = \sum_{i=1}^{n} c_i\vv_i
\]
\end{definition}

\begin{intuition}
``线性 (linear)''这个名字从何而来? 如果固定其中一个标量, 让另一个标量变化, 所得向量的箭头终点会沿一条\textbf{直线}移动。

例如, 若固定 $b = 1$, 让 $a$ 变化, 则向量 $a\vv + \vw$ 会沿一条与 $\vv$ 平行的直线移动。
\end{intuition}

\begin{example}
设 $\vv = \twovec{1}{2}$, $\vw = \twovec{3}{-1}$。一些线性组合如下:
\begin{align*}
2\vv + 1\vw &= 2\twovec{1}{2} + 1\twovec{3}{-1} = \twovec{2}{4} + \twovec{3}{-1} = \twovec{5}{3} \\
-1\vv + 3\vw &= -\twovec{1}{2} + 3\twovec{3}{-1} = \twovec{-1}{-2} + \twovec{9}{-3} = \twovec{8}{-5} \\
0\vv + 0\vw &= \twovec{0}{0} = \vzero
\end{align*}
\end{example}

% ============================================
\section{张成}
% ============================================

\begin{definition}[张成 (span)]
一组向量的\vocab{张成}就是这些向量所有可能的线性组合构成的集合:
\[
\spn\{\vv_1, \vv_2, \ldots, \vv_n\} = \{c_1\vv_1 + c_2\vv_2 + \cdots + c_n\vv_n \mid c_i \in \R\}
\]
\end{definition}

\begin{intuition}
关键问题是: 给定一些向量, 只用两种基本运算 (加法与标量乘法), 我们能到达哪些向量?

答案就是: 这些向量的\textbf{张成}。
\end{intuition}

% -------------------------------------------
\subsection{二维空间中的张成}
% -------------------------------------------

\begin{theorem}
对于 $\Rtwo$ 中的两个向量 $\vv$ 与 $\vw$, 有以下三种可能情形:
\begin{enumerate}
    \item 若 $\vv$ 与 $\vw$ 不共线: $\spn\{\vv, \vw\} = \Rtwo$ (整个平面)
    \item 若 $\vv$ 与 $\vw$ 共线 (但非零): $\spn\{\vv, \vw\}$ 是一条直线
    \item 若两者均为零向量: $\spn\{\vv, \vw\} = \{\vzero\}$ (只有原点)
\end{enumerate}
\end{theorem}

\begin{center}
\begin{tikzpicture}[scale=0.9]
    % 情形 1: 整个平面
    \begin{scope}[shift={(0,0)}]
        \fill[blue!10] (-2,-2) rectangle (2,2);
        \draw[axis] (-2,0) -- (2,0);
        \draw[axis] (0,-2) -- (0,2);
        \draw[vector, red!70, very thick] (0,0) -- (1,0.5) node[right] {$\vv$};
        \draw[vector, green!60!black, very thick] (0,0) -- (-0.3,1) node[above] {$\vw$};
        \node at (0,-2.5) {情形 1: 整个平面};
    \end{scope}

    % 情形 2: 一条直线
    \begin{scope}[shift={(6,0)}]
        \draw[blue!30, very thick] (-2,-1) -- (2,1);
        \draw[axis] (-2,0) -- (2,0);
        \draw[axis] (0,-2) -- (0,2);
        \draw[vector, red!70, very thick] (0,0) -- (1,0.5) node[right] {$\vv$};
        \draw[vector, green!60!black, very thick] (0,0) -- (1.6,0.8) node[above] {$\vw$};
        \node at (0,-2.5) {情形 2: 一条直线};
    \end{scope}

    % 情形 3: 只有原点
    \begin{scope}[shift={(12,0)}]
        \draw[axis] (-2,0) -- (2,0);
        \draw[axis] (0,-2) -- (0,2);
        \fill[blue!70] (0,0) circle (3pt);
        \node at (0,-2.5) {情形 3: 只有原点};
    \end{scope}
\end{tikzpicture}
\end{center}

% -------------------------------------------
\subsection{三维空间中的张成}
% -------------------------------------------

\begin{theorem}
对于 $\Rthree$ 中的向量:
\begin{itemize}
    \item \textbf{一个非零向量:} 张成是一条\textbf{直线}
    \item \textbf{两个不共线的向量:} 张成是一个\textbf{平面}
    \item \textbf{三个不共面的向量:} 张成是整个 $\Rthree$
\end{itemize}
\end{theorem}

\begin{intuition}
想象三维空间中的两个向量。它们的线性组合构成一个过原点的平面。现在再加入一个不在该平面上的第三个向量, 就好比把这个平面在空间中``扫动''起来, 从而覆盖整个三维空间。
\end{intuition}

% ============================================
\section{线性无关与线性相关}
% ============================================

\begin{definition}[线性相关 (linear dependence)]
若一组向量中可以去掉其中一个而不改变张成, 则称这组向量\vocab{线性相关}, 换句话说, 其中至少有一个向量是``多余的''。

数学表述: 向量 $\vv_1, \ldots, \vv_n$ 线性相关, 当且仅当存在不全为零的标量 $c_1, \ldots, c_n$, 使得:
\[
c_1\vv_1 + c_2\vv_2 + \cdots + c_n\vv_n = \vzero
\]
\end{definition}

\begin{definition}[线性无关 (linear independence)]
若一组向量中没有一个是多余的---即每个向量都为张成增加了新的维度, 则称这组向量\vocab{线性无关}。

数学表述: 向量 $\vv_1, \ldots, \vv_n$ 线性无关, 当且仅当:
\[
c_1\vv_1 + c_2\vv_2 + \cdots + c_n\vv_n = \vzero \implies c_1 = c_2 = \cdots = c_n = 0
\]
\end{definition}

\begin{example}
考虑向量 $\vv = \twovec{2}{1}$ 与 $\vw = \twovec{4}{2}$。

这两个向量\textbf{线性相关}, 因为 $\vw = 2\vv$。我们可以写成:
\[
2\vv - \vw = \vzero \quad \text{即} \quad 2\twovec{2}{1} - \twovec{4}{2} = \twovec{0}{0}
\]
\end{example}

\begin{example}
向量 $\vv = \twovec{1}{0}$ 与 $\vw = \twovec{0}{1}$ 线性无关。

若 $a\vv + b\vw = \vzero$:
\[
a\twovec{1}{0} + b\twovec{0}{1} = \twovec{a}{b} = \twovec{0}{0}
\]
于是 $a = 0$ 且 $b = 0$。只有当所有系数都为零时, 才能得到零向量。
\end{example}

\begin{warning}
在 $n$ 维空间中, 至多有 $n$ 个向量能线性无关。如果向量个数超过 $n$, 它们一定线性相关。
\end{warning}

% ============================================
\section{基}
% ============================================

\begin{definition}[基 (basis)]
向量空间的一组\vocab{基}是满足以下两个条件的一组向量:
\begin{enumerate}
    \item \textbf{线性无关}
    \item \textbf{张成}整个空间
\end{enumerate}
\end{definition}

\begin{theorem}
$\Rtwo$ 的标准基为:
\[
\left\{ \vi = \twovec{1}{0}, \; \vj = \twovec{0}{1} \right\}
\]
$\Rthree$ 的标准基为:
\[
\left\{ \vi = \threevec{1}{0}{0}, \; \vj = \threevec{0}{1}{0}, \; \vk = \threevec{0}{0}{1} \right\}
\]
\end{theorem}

\begin{intuition}
基就像描述向量的一门``语言''。当我们说 $\twovec{3}{2}$ 时, 实际是在说``取 3 个第一个基向量, 再取 2 个第二个基向量''。

如果换一组基, 同一个向量的坐标也会不同---就像把一句话翻译成另一种语言。
\end{intuition}

% -------------------------------------------
\subsection{非标准基}
% -------------------------------------------

\begin{example}
以下集合也是 $\Rtwo$ 的一组基:
\[
\mathcal{B} = \left\{ \vb_1 = \twovec{1}{1}, \; \vb_2 = \twovec{1}{-1} \right\}
\]

\textbf{线性无关性的证明:} 若 $a\vb_1 + b\vb_2 = \vzero$:
\[
a\twovec{1}{1} + b\twovec{1}{-1} = \twovec{a+b}{a-b} = \twovec{0}{0}
\]
于是 $a + b = 0$ 且 $a - b = 0$, 解得 $a = b = 0$。

\textbf{张成整个平面:} 因为 $\Rtwo$ 中有两个线性无关的向量, 它们张成整个平面。
\end{example}

\begin{center}
\begin{tikzpicture}[scale=1.5]
    % 网格
    \draw[grid] (-2,-2) grid (2,2);

    % 坐标轴
    \draw[axis] (-2,0) -- (2,0) node[right] {$x$};
    \draw[axis] (0,-2) -- (0,2) node[above] {$y$};

    % 标准基
    \draw[basis, -{Stealth[length=3mm]}] (0,0) -- (1,0) node[below right] {$\vi$};
    \draw[basis, -{Stealth[length=3mm]}] (0,0) -- (0,1) node[above left] {$\vj$};

    % 非标准基
    \draw[vector, blue!70, very thick] (0,0) -- (1,1) node[above right] {$\vb_1$};
    \draw[vector, green!60!black, very thick] (0,0) -- (1,-1) node[below right] {$\vb_2$};
\end{tikzpicture}
\end{center}

% ============================================
\section{把向量看作点}
% ============================================

\begin{remark}
有时把向量看作一支箭头不如把它看作一个\textbf{点}---即向量箭头终点落到的那个点, 会更加方便。

\begin{itemize}
    \item 当思考单个具体的向量时: 把它想象成一支\textbf{箭头}
    \item 当思考一族向量时: 把它们想象成一些\textbf{点}
\end{itemize}
\end{remark}

\begin{center}
\begin{tikzpicture}[scale=1]
    % 左图: 作为箭头的向量
    \begin{scope}[shift={(0,0)}]
        \draw[axis] (-0.5,0) -- (3,0);
        \draw[axis] (0,-0.5) -- (0,3);
        \draw[vector, blue!70] (0,0) -- (1,2);
        \draw[vector, red!70] (0,0) -- (2,1);
        \draw[vector, green!60!black] (0,0) -- (2.5,2.5);
        \node at (1.5,-1) {作为箭头的向量};
    \end{scope}

    % 右图: 作为点的向量
    \begin{scope}[shift={(6,0)}]
        \draw[axis] (-0.5,0) -- (3,0);
        \draw[axis] (0,-0.5) -- (0,3);
        \fill[blue!70] (1,2) circle (3pt);
        \fill[red!70] (2,1) circle (3pt);
        \fill[green!60!black] (2.5,2.5) circle (3pt);
        \node at (1.5,-1) {作为点的向量};
    \end{scope}
\end{tikzpicture}
\end{center}

% ============================================
\section{核心概念总结}
% ============================================

\begin{summary}
\begin{description}
    \item[线性组合:] $c_1\vv_1 + c_2\vv_2 + \cdots + c_n\vv_n$, 其中系数为任意标量
    \item[张成:] 所有可能的线性组合构成的集合
    \item[线性无关:] 没有向量是多余的
    \item[线性相关:] 至少有一个向量可以去掉
    \item[基:] 线性无关且张成整个空间的向量组
\end{description}
\end{summary}

% ============================================
\section{练习}
% ============================================

\begin{exercise}
把向量 $\vv = \twovec{4}{6}$ 写成 $\vi$ 与 $\vj$ 的线性组合。
\end{exercise}

\begin{exercise}
以下向量是否线性无关? 请说明理由。
\[
\va = \twovec{1}{3}, \quad \vb = \twovec{2}{6}
\]
\end{exercise}

\begin{exercise}
以下向量是否线性无关?
\[
\vv_1 = \twovec{1}{0}, \quad \vv_2 = \twovec{1}{1}, \quad \vv_3 = \twovec{0}{1}
\]
\textit{提示: 在 $\Rtwo$ 中, 最多有几个线性无关的向量?}
\end{exercise}

\begin{exercise}
描述以下向量的张成:
\[
\vv = \threevec{1}{0}{0}, \quad \vw = \threevec{0}{1}{0}
\]
\end{exercise}

\begin{exercise}[挑战]
证明以下集合是 $\Rtwo$ 的一组基:
\[
\left\{ \twovec{2}{1}, \twovec{1}{3} \right\}
\]
然后求向量 $\twovec{5}{7}$ 在这组新基下的坐标。
\end{exercise}

\begin{exercise}
给定 $\Rthree$ 中的三个向量:
\[
\vv_1 = \threevec{1}{0}{1}, \quad \vv_2 = \threevec{0}{1}{1}, \quad \vv_3 = \threevec{1}{1}{2}
\]
这些向量是否线性无关? 它们的张成是什么?
\end{exercise}

\begin{problem}
证明: 若 $\{\vv_1, \vv_2\}$ 是 $\Rtwo$ 的一组基, 则 $\{2\vv_1, 3\vv_2\}$ 也是一组基。
\end{problem}

\begin{problem}
在 $\Rtwo$ 中能否找到三个线性无关的向量? 为什么?
\end{problem}


% 第 3 章: 线性变换与矩阵
% lecture03.tex - 线性变换与矩阵
% 第 3 章: 线性变换与矩阵

\chapter{线性变换与矩阵}
\label{ch:transformations}

\begin{abstract}
线性变换 (linear transformation) 是线性代数的核心。在本章中, 我们将学习任何线性变换都可以用一个矩阵 (matrix) 来表示, 以及矩阵乘法与变换复合之间的关系。
\end{abstract}

% ============================================
\section{什么是变换?}
% ============================================

\begin{definition}[变换 (transformation)]
\vocab{变换}是一个把向量变成另一个向量的函数:
\[
T: \Rtwo \to \Rtwo
\]
每个输入向量都被映射到一个输出向量。
\end{definition}

\begin{intuition}
我们用``变换''而不用``函数'', 是因为我们想强调\textbf{运动}。想象空间中的每个向量都被``移动''到了新的位置。

再想象坐标系的网格线。变换会拉伸、压缩、旋转这些网格, 或以其他方式改变它们。
\end{intuition}

% ============================================
\section{线性变换}
% ============================================

\begin{definition}[线性变换 (linear transformation)]
变换 $T$ 是\vocab{线性变换}, 当且仅当以下两个条件成立:
\begin{enumerate}
    \item \textbf{可加性:} $T(\vv + \vw) = T(\vv) + T(\vw)$
    \item \textbf{齐次性:} $T(c\vv) = c \cdot T(\vv)$
\end{enumerate}
其中 $\vv, \vw$ 为任意向量, $c$ 为任意标量。
\end{definition}

\begin{intuition}
线性变换可以从两个几何性质来识别:
\begin{enumerate}
    \item \textbf{直线变换后仍是直线} (不会变弯)
    \item \textbf{原点保持不动}
\end{enumerate}

如果变换后网格线仍然平行且等距, 那它就是线性的。
\end{intuition}

\begin{center}
\begin{tikzpicture}[scale=0.8]
    % 原网格
    \begin{scope}[shift={(0,0)}]
        \draw[grid] (-2,-2) grid (2,2);
        \draw[axis] (-2,0) -- (2,0);
        \draw[axis] (0,-2) -- (0,2);
        \draw[basis, -{Stealth[length=3mm]}] (0,0) -- (1,0) node[below] {$\vi$};
        \draw[basis, -{Stealth[length=3mm]}] (0,0) -- (0,1) node[left] {$\vj$};
        \node at (0,-3) {变换前};
    \end{scope}

    % 箭头
    \draw[-{Stealth[length=5mm]}, very thick] (3,0) -- (5,0) node[midway, above] {$T$};

    % 变换后的网格 (错切)
    \begin{scope}[shift={(8,0)}]
        \draw[blue!30, thin] (-2,-2) -- (0,2) -- (4,2) -- (2,-2) -- cycle;
        \foreach \y in {-1,0,1} {
            \draw[blue!30, thin] ({-2+\y},-2) -- ({\y+2},2);
        }
        \foreach \x in {-2,-1,0,1,2} {
            \draw[blue!30, thin] ({\x-1},-2) -- ({\x+1},2);
        }
        \draw[axis] (-2,0) -- (4,0);
        \draw[axis] (0,-2) -- (2,2);
        \draw[transformed, -{Stealth[length=3mm]}, very thick] (0,0) -- (1,0) node[below] {$T(\vi)$};
        \draw[transformed, -{Stealth[length=3mm]}, very thick] (0,0) -- (1,1) node[above left] {$T(\vj)$};
        \node at (1,-3) {变换后};
    \end{scope}
\end{tikzpicture}
\end{center}

% ============================================
\section{线性变换的矩阵}
% ============================================

\begin{theorem}[基本定理]
每个线性变换 $T: \Rtwo \to \Rtwo$ 完全由 $T$ 把基向量 (basis) $\vi$ 与 $\vj$ 送到哪里决定。
\end{theorem}

\begin{proof}
任何向量 $\vv = \twovec{x}{y}$ 都可以写成: $\vv = x\vi + y\vj$

利用线性性质:
\[
T(\vv) = T(x\vi + y\vj) = xT(\vi) + yT(\vj)
\]
因此只要知道 $T(\vi)$ 与 $T(\vj)$, 就能算出任何 $\vv$ 的 $T(\vv)$。
\end{proof}

\begin{definition}[变换矩阵 (transformation matrix)]
若 $T(\vi) = \twovec{a}{c}$, $T(\vj) = \twovec{b}{d}$, 则 $T$ 的\vocab{变换矩阵}为:
\[
\mA = \twomat{a}{b}{c}{d}
\]
第一列是 $\vi$ 的落点, 第二列是 $\vj$ 的落点。
\end{definition}

\begin{important}
\textbf{黄金法则:} 矩阵的列向量 (column vector) = 基向量的落点
\end{important}

% ============================================
\section{矩阵与向量相乘}
% ============================================

\begin{definition}[矩阵乘以向量]
\[
\twomat{a}{b}{c}{d} \twovec{x}{y} = x\twovec{a}{c} + y\twovec{b}{d} = \twovec{ax + by}{cx + dy}
\]
\end{definition}

\begin{intuition}
矩阵乘以向量意味着:
\begin{enumerate}
    \item 把输入的 $x$ 分量乘以第一列
    \item 把输入的 $y$ 分量乘以第二列
    \item 把结果加起来
\end{enumerate}
这正是新基向量的线性组合!
\end{intuition}

\begin{example}
逆时针旋转 $90°$:
\[
T(\vi) = \twovec{0}{1}, \quad T(\vj) = \twovec{-1}{0}
\]
于是旋转矩阵为:
\[
\mR = \twomat{0}{-1}{1}{0}
\]
验证: $\mR\twovec{1}{0} = \twovec{0}{1}$ \checkmark
\end{example}

% ============================================
\section{重要变换举例}
% ============================================

\subsection{旋转}

\begin{definition}[旋转矩阵]
逆时针旋转 $\theta$ 角:
\[
\mR_\theta = \twomat{\cos\theta}{-\sin\theta}{\sin\theta}{\cos\theta}
\]
\end{definition}

\begin{example}
旋转 $45°$:
\[
\mR_{45°} = \twomat{\frac{\sqrt{2}}{2}}{-\frac{\sqrt{2}}{2}}{\frac{\sqrt{2}}{2}}{\frac{\sqrt{2}}{2}}
\]
\end{example}

\subsection{缩放}

\begin{definition}[缩放矩阵]
沿 $x$ 轴缩放 $s_x$ 倍、沿 $y$ 轴缩放 $s_y$ 倍:
\[
\mS = \twomat{s_x}{0}{0}{s_y}
\]
\end{definition}

\begin{practical}
\textbf{在计算机图形学中的应用:} 要放大或缩小一幅图像:
\begin{itemize}
    \item $s_x = s_y = 2$: 图像变为原来的两倍
    \item $s_x = 1, s_y = 0.5$: 图像在竖直方向被压缩
\end{itemize}
\end{practical}

\subsection{错切}

\begin{definition}[水平错切矩阵]
\[
\mH = \twomat{1}{k}{0}{1}
\]
其中 $k$ 为错切量。
\end{definition}

\begin{center}
\begin{tikzpicture}[scale=1]
    % 原图形
    \begin{scope}[shift={(0,0)}]
        \draw[fill=blue!20] (0,0) -- (1,0) -- (1,1) -- (0,1) -- cycle;
        \draw[basis] (0,0) -- (1,0) node[midway, below] {$\vi$};
        \draw[basis] (0,0) -- (0,1) node[midway, left] {$\vj$};
        \node at (0.5,-0.7) {原图形};
    \end{scope}

    % 错切后
    \begin{scope}[shift={(4,0)}]
        \draw[fill=blue!20] (0,0) -- (1,0) -- (2,1) -- (1,1) -- cycle;
        \draw[transformed, very thick] (0,0) -- (1,0) node[midway, below] {$T(\vi)$};
        \draw[transformed, very thick] (0,0) -- (1,1) node[midway, left] {$T(\vj)$};
        \node at (1,-0.7) {错切 ($k=1$)};
    \end{scope}
\end{tikzpicture}
\end{center}

\subsection{反射}

\begin{example}
关于 $x$ 轴的反射:
\[
\twomat{1}{0}{0}{-1}
\]

关于直线 $y = x$ 的反射:
\[
\twomat{0}{1}{1}{0}
\]
\end{example}

% ============================================
\section{矩阵乘法即变换的复合}
% ============================================

\begin{theorem}[变换的复合]
若 $\mA$ 是变换 $T_1$ 的矩阵, $\mB$ 是变换 $T_2$ 的矩阵, 则:
\[
\mB \cdot \mA = \text{矩阵 } (T_2 \circ T_1)
\]
即先做 $T_1$, 再做 $T_2$。
\end{theorem}

\begin{warning}
顺序很重要! 一般来说 $\mA\mB \neq \mB\mA$。

先旋转再错切 $\neq$ 先错切再旋转
\end{warning}

\begin{definition}[矩阵乘法 (matrix multiplication)]
\[
\twomat{a}{b}{c}{d} \twomat{e}{f}{g}{h} = \twomat{ae+bg}{af+bh}{ce+dg}{cf+dh}
\]
\end{definition}

\begin{intuition}
计算 $\mB\mA$ 的列:
\begin{itemize}
    \item 第一列: $\mB$ 乘以 $\mA$ 的第一列 = $\vi$ 经过两次变换后的落点
    \item 第二列: $\mB$ 乘以 $\mA$ 的第二列 = $\vj$ 经过两次变换后的落点
\end{itemize}
\end{intuition}

\begin{example}
旋转 $90°$ 后再错切:
\[
\underbrace{\twomat{1}{1}{0}{1}}_{\text{错切}} \underbrace{\twomat{0}{-1}{1}{0}}_{\text{旋转}} = \twomat{1}{-1}{1}{0}
\]

验证: $\vi \to \twovec{0}{1} \to \twovec{1}{1}$ \checkmark
\end{example}

% ============================================
\section{实际应用}
% ============================================

\begin{practical}
\textbf{计算机图形学与电子游戏}

游戏中的每个物体都由一组向量 (顶点) 描述。要移动、旋转或缩放一个物体, 只需把所有顶点乘以相应的矩阵。

一段动画 = 一连串矩阵乘法
\end{practical}

\begin{practical}
\textbf{图像处理}

模糊、锐化、边缘检测等图像滤波器都是用矩阵乘法实现的。
\end{practical}

% ============================================
\section{练习}
% ============================================

\begin{exercise}
写出把 $\vi$ 送到 $\twovec{2}{1}$、把 $\vj$ 送到 $\twovec{-1}{3}$ 的变换的矩阵。
\end{exercise}

\begin{exercise}
计算以下乘积:
\[
\twomat{1}{2}{3}{4} \twovec{5}{6}
\]
\end{exercise}

\begin{exercise}
写出旋转 $180°$ 的矩阵, 并证明它等于 $-\mI$。
\end{exercise}

\begin{exercise}
把以下两个矩阵按两种顺序相乘, 并证明 $\mA\mB \neq \mB\mA$:
\[
\mA = \twomat{1}{2}{0}{1}, \quad \mB = \twomat{0}{1}{1}{0}
\]
\end{exercise}

\begin{exercise}[挑战]
求关于直线 $y = 2x$ 的反射矩阵。

\textit{提示: 这条直线与 $x$ 轴的夹角为 $\arctan(2)$。}
\end{exercise}

\begin{problem}
证明: 旋转 $\alpha$ 角后再旋转 $\beta$ 角, 等于旋转 $\alpha + \beta$ 角。
\end{problem}

\begin{problem}
若 $\mA^2 = \mA$ (即 $\mA$ 为投影矩阵), 它的几何意义是什么?
\end{problem}


% 第 4 章: 三维变换与行列式
% lecture04.tex - 三维变换与行列式
% 第 4 章: 三维变换与行列式 (3D Transformations and Determinants)

\chapter{三维变换与行列式}
\label{ch:determinant}

\begin{abstract}
行列式 (determinant) 是一个数, 它度量了线性变换将空间``拉伸''或``压缩''了多少。在本章中, 我们将学习行列式的几何意义以及如何计算行列式。
\end{abstract}

% ============================================
\section{三维空间中的变换}
% ============================================

\begin{definition}[三维变换矩阵]
$\Rthree$ 中的线性变换可以用一个 $3 \times 3$ 矩阵表示:
\[
\mA = \threemat{a_{11}}{a_{12}}{a_{13}}{a_{21}}{a_{22}}{a_{23}}{a_{31}}{a_{32}}{a_{33}}
\]
它的各列分别是 $\vi$、$\vj$ 和 $\vk$ 落到的位置。
\end{definition}

\begin{intuition}
二维中的同样道理在三维中也适用。想象一张三维网格被拉伸、压缩、旋转或错切。直线仍然保持为直线, 原点保持不动。
\end{intuition}

% ============================================
\section{行列式的几何意义}
% ============================================

\begin{definition}[行列式---几何定义]
矩阵的\vocab{行列式 (determinant)}是相应的变换缩放面积 (二维) 或体积 (三维) 的倍数。

若 $\mA$ 是变换 $T$ 的矩阵:
\[
\det(\mA) = \frac{\text{变换后的面积/体积}}{\text{变换前的面积/体积}}
\]
\end{definition}

\begin{center}
\begin{tikzpicture}[scale=1.2]
    % 原始单位正方形
    \begin{scope}[shift={(0,0)}]
        \fill[blue!20] (0,0) -- (1,0) -- (1,1) -- (0,1) -- cycle;
        \draw[thick] (0,0) -- (1,0) -- (1,1) -- (0,1) -- cycle;
        \draw[basis] (0,0) -- (1,0);
        \draw[basis] (0,0) -- (0,1);
        \node at (0.5,0.5) {$1$};
        \node at (0.5,-0.5) {面积 = $1$};
    \end{scope}

    \draw[-{Stealth}, very thick] (1.5,0.5) -- (2.5,0.5) node[midway, above] {$T$};

    % 变换后的平行四边形
    \begin{scope}[shift={(3,0)}]
        \fill[green!20] (0,0) -- (2,0) -- (3,1.5) -- (1,1.5) -- cycle;
        \draw[thick] (0,0) -- (2,0) -- (3,1.5) -- (1,1.5) -- cycle;
        \draw[transformed, very thick] (0,0) -- (2,0);
        \draw[transformed, very thick] (0,0) -- (1,1.5);
        \node at (1.5,0.75) {$3$};
        \node at (1.5,-0.5) {面积 = $3$};
    \end{scope}
\end{tikzpicture}
\end{center}

\begin{important}
行列式就是面积/体积的缩放因子。

若 $\det(\mA) = 3$, 则每个图形经过变换后都变为原来的 $3$ 倍大。
\end{important}

% ============================================
\section{行列式的符号}
% ============================================

\begin{theorem}[符号的意义]
\begin{itemize}
    \item $\det(\mA) > 0$: 空间的定向保持不变
    \item $\det(\mA) < 0$: 空间的定向被反转 (像镜子一样)
    \item $\det(\mA) = 0$: 空间被压缩到更低的维度
\end{itemize}
\end{theorem}

\begin{intuition}
在二维中, 若 $\vj$ 在 $\vi$ 的左侧, 则定向为``正''。若某个变换反转了这种关系 (比如反射), 行列式就变为负数。

在三维中, 用右手定则判断: 若手指从 $\vi$ 向 $\vj$ 卷曲时拇指指向 $\vk$, 则定向为正。
\end{intuition}

% ============================================
\section{行列式的计算}
% ============================================

\subsection{二阶矩阵的行列式}

\begin{definition}
\[
\det\twomat{a}{b}{c}{d} = ad - bc
\]
\end{definition}

\begin{example}
\[
\det\twomat{3}{1}{0}{2} = 3 \times 2 - 1 \times 0 = 6
\]
每个图形的面积都被乘以 $6$。
\end{example}

\subsection{三阶矩阵的行列式}

\begin{definition}[展开公式]
\[
\det\threemat{a}{b}{c}{d}{e}{f}{g}{h}{i} = a(ei-fh) - b(di-fg) + c(dh-eg)
\]
\end{definition}

\begin{intuition}
$3 \times 3$ 行列式等于由矩阵的三个列向量张成的平行六面体的体积 (带符号)。
\end{intuition}

\begin{example}
\[
\det\threemat{1}{2}{3}{4}{5}{6}{7}{8}{9}
\]
\begin{align*}
&= 1(5 \times 9 - 6 \times 8) - 2(4 \times 9 - 6 \times 7) + 3(4 \times 8 - 5 \times 7) \\
&= 1(45-48) - 2(36-42) + 3(32-35) \\
&= 1(-3) - 2(-6) + 3(-3) \\
&= -3 + 12 - 9 = 0
\end{align*}
行列式为零说明这三个列向量落在同一个平面上!
\end{example}

% ============================================
\section{行列式的性质}
% ============================================

\begin{theorem}[主要性质]
\begin{enumerate}
    \item $\det(\mI) = 1$
    \item $\det(\mA\mB) = \det(\mA) \cdot \det(\mB)$
    \item $\det(\mA\trans) = \det(\mA)$
    \item 对 $n \times n$ 矩阵有 $\det(c\mA) = c^n \det(\mA)$
    \item 若某一行/列全为零: $\det(\mA) = 0$
    \item 若两行/两列相等: $\det(\mA) = 0$
\end{enumerate}
\end{theorem}

\begin{intuition}
性质 $\det(\mA\mB) = \det(\mA) \cdot \det(\mB)$ 非常重要:

若 $\mA$ 把面积放大 $3$ 倍, $\mB$ 把面积放大 $2$ 倍, 则它们的复合把面积放大 $3 \times 2 = 6$ 倍。
\end{intuition}

% ============================================
\section{零行列式与空间坍缩}
% ============================================

\begin{theorem}
$\det(\mA) = 0$ 当且仅当:
\begin{itemize}
    \item $\mA$ 的各列线性相关
    \item 相应的变换把空间压缩到更低的维度
\end{itemize}
\end{theorem}

\begin{center}
\begin{tikzpicture}[scale=1]
    % 二维到一维
    \begin{scope}[shift={(0,0)}]
        \fill[blue!20] (0,0) -- (1,0) -- (1,1) -- (0,1) -- cycle;
        \node at (0.5,-0.5) {$\Rtwo$};
    \end{scope}

    \draw[-{Stealth}, very thick] (1.5,0.5) -- (3,0.5) node[midway, above] {$\det = 0$};

    \begin{scope}[shift={(3.5,0)}]
        \draw[blue!70, ultra thick] (-0.5,0.5) -- (2,0.5);
        \node at (0.75,-0.5) {一条直线};
    \end{scope}
\end{tikzpicture}
\end{center}

\begin{practical}
\textbf{检验线性无关性:}

想知道空间中的三个向量是否线性无关? 把它们作为列排成一个矩阵。若 $\det \neq 0$, 则它们线性无关。
\end{practical}

% ============================================
\section{实际应用}
% ============================================

\begin{practical}
\textbf{由顶点计算三角形面积}

若三角形的顶点为 $(x_1, y_1)$、$(x_2, y_2)$、$(x_3, y_3)$:
\[
\text{Area} = \frac{1}{2} \left| \det\threemat{x_1}{y_1}{1}{x_2}{y_2}{1}{x_3}{y_3}{1} \right|
\]
\end{practical}

\begin{practical}
\textbf{物理: 力矩与力}

行列式用于计算叉积, 而叉积出现在力矩、角动量和电磁场的计算中。
\end{practical}

% ============================================
\section{练习}
% ============================================

\begin{exercise}
计算以下矩阵的行列式:
\begin{enumerate}[label=(\alph*)]
    \item $\twomat{3}{7}{1}{4}$
    \item $\twomat{2}{6}{1}{3}$
    \item $\twomat{-1}{2}{3}{-6}$
\end{enumerate}
\end{exercise}

\begin{exercise}
若 $\det(\mA) = 5$ 且 $\det(\mB) = -2$, 则 $\det(\mA\mB)$ 等于多少?
\end{exercise}

\begin{exercise}
计算顶点为 $(0,0)$、$(3,1)$、$(1,4)$、$(4,5)$ 的平行四边形的面积。
\end{exercise}

\begin{exercise}
计算旋转角为 $\theta$ 的旋转矩阵的行列式, 并解释结果的含义。
\end{exercise}

\begin{exercise}[挑战]
证明对任何矩阵 $\mA$:
\[
\det(\mA\inv) = \frac{1}{\det(\mA)}
\]
(假设 $\mA$ 可逆)
\end{exercise}

\begin{problem}
计算由以下向量张成的平行六面体的体积:
\[
\va = \threevec{1}{0}{2}, \quad \vb = \threevec{0}{3}{1}, \quad \vc = \threevec{2}{1}{0}
\]
\end{problem}


% 第 5 章: 逆矩阵、列空间与零空间
% lecture05.tex - 逆矩阵、列空间与零空间
% 第 5 章: 逆矩阵、列空间与零空间 (Inverse Matrices, Column Space, and Null Space)

\chapter{逆矩阵、列空间与零空间}
\label{ch:inverse}

\begin{abstract}
在本章中, 我们将学习逆矩阵 (inverse matrix)、线性方程组的求解, 以及与矩阵相关的重要空间 (列空间、零空间和秩)。
\end{abstract}

% ============================================
\section{线性方程组}
% ============================================

\begin{definition}[方程组的矩阵形式]
线性方程组:
\[
\system{
a_{11}x_1 + a_{12}x_2 + \cdots + a_{1n}x_n &= b_1 \\
a_{21}x_1 + a_{22}x_2 + \cdots + a_{2n}x_n &= b_2 \\
&\vdots \\
a_{m1}x_1 + a_{m2}x_2 + \cdots + a_{mn}x_n &= b_m
}
\]
可以写成 $\mA\vx = \vb$。
\end{definition}

\begin{intuition}
把方程 $\mA\vx = \vb$ 理解为:
\begin{center}
``哪个向量 $\vx$ 经过变换 $\mA$ 后会落到 $\vb$ 上?''
\end{center}
换句话说: 把变换 $\mA$``倒回去'', 从 $\vb$ 回到 $\vx$。
\end{intuition}

% ============================================
\section{逆矩阵}
% ============================================

\begin{definition}[逆矩阵]
若矩阵 $\mA\inv$ 满足:
\[
\mA\inv \mA = \mA \mA\inv = \mI
\]
则称它是矩阵 $\mA$ 的\vocab{逆 (inverse)}, 其中 $\mI$ 是单位矩阵。
\end{definition}

\begin{intuition}
若 $\mA$ 是一个变换, 则 $\mA\inv$ 就是抵消 $\mA$ 效果的变换:
\begin{itemize}
    \item 若 $\mA$ 是旋转 $90°$, 则 $\mA\inv$ 是旋转 $-90°$
    \item 若 $\mA$ 是放大 $2$ 倍, 则 $\mA\inv$ 是缩小为 $\frac{1}{2}$
\end{itemize}
\end{intuition}

\begin{theorem}[用逆矩阵解方程组]
若 $\mA$ 可逆, 则 $\mA\vx = \vb$ 的唯一解为:
\[
\vx = \mA\inv \vb
\]
\end{theorem}

\subsection{二阶矩阵的求逆公式}

\begin{theorem}
若 $\mA = \twomat{a}{b}{c}{d}$ 且 $\det(\mA) \neq 0$:
\[
\mA\inv = \frac{1}{ad-bc} \twomat{d}{-b}{-c}{a}
\]
\end{theorem}

\begin{example}
\[
\mA = \twomat{3}{1}{2}{4}, \quad \det(\mA) = 12 - 2 = 10
\]
\[
\mA\inv = \frac{1}{10} \twomat{4}{-1}{-2}{3} = \twomat{0.4}{-0.1}{-0.2}{0.3}
\]
\end{example}

% ============================================
\section{逆何时存在}
% ============================================

\begin{theorem}
矩阵 $\mA$ 可逆当且仅当 $\det(\mA) \neq 0$。
\end{theorem}

\begin{intuition}
若 $\det(\mA) = 0$, 则变换 $\mA$ 会压缩空间 (例如把平面压成直线)。这样的变换不可逆---信息已经丢失了。

这就像把几个数加起来: 只看总和, 无法恢复原来的每个数。
\end{intuition}

% ============================================
\section{列空间}
% ============================================

\begin{definition}[列空间]
矩阵 $\mA$ 的\vocab{列空间 (column space)}, 记为 $\Col(\mA)$, 是 $\mA$ 的各列张成的空间:
\[
\Col(\mA) = \spn\{\text{columns of } \mA\}
\]
\end{definition}

\begin{intuition}
列空间就是变换 $\mA$ 的所有可能输出组成的集合。

问题``$\mA\vx = \vb$ 是否有解?''等价于``$\vb$ 是否在 $\mA$ 的列空间中?''
\end{intuition}

\begin{example}
\[
\mA = \twomat{1}{3}{2}{6}
\]
第二列是第一列的 $3$ 倍, 所以:
\[
\Col(\mA) = \spn\left\{\twovec{1}{2}\right\} = \text{一条过原点的直线}
\]
\end{example}

% ============================================
\section{零空间}
% ============================================

\begin{definition}[零空间]
矩阵 $\mA$ 的\vocab{零空间 (null space)} (或核) 是所有被 $\mA$ 送到零向量的向量组成的集合:
\[
\Null(\mA) = \{\vx \mid \mA\vx = \vzero\}
\]
\end{definition}

\begin{intuition}
零空间就是被变换 $\mA$``压扁''到零的那些向量。

若 $\det(\mA) \neq 0$: 只有零向量被压扁, 所以 $\Null(\mA) = \{\vzero\}$。

若 $\det(\mA) = 0$: 整条直线或整个平面都被压缩到原点。
\end{intuition}

\begin{example}
对矩阵 $\mA = \twomat{1}{2}{2}{4}$:

解 $\mA\vx = \vzero$:
\[
\twomat{1}{2}{2}{4}\twovec{x}{y} = \twovec{0}{0}
\]
方程为: $x + 2y = 0$, 所以 $x = -2y$。

零空间为: $\Null(\mA) = \left\{ t\twovec{-2}{1} \mid t \in \R \right\}$ (一条直线)
\end{example}

% ============================================
\section{秩}
% ============================================

\begin{definition}[秩]
矩阵 $\mA$ 的\vocab{秩 (rank)}等于:
\begin{itemize}
    \item 列空间的维数
    \item 线性无关的列的个数
    \item 线性无关的行的个数
\end{itemize}
\[
\rank(\mA) = \dim(\Col(\mA))
\]
\end{definition}

\begin{theorem}[秩---零化度定理]
对 $m \times n$ 矩阵:
\[
\rank(\mA) + \dim(\Null(\mA)) = n
\]
(即列的个数)
\end{theorem}

\begin{intuition}
秩就是变换保留下来的维数。

$\dim(\Null(\mA))$ 就是丢失掉的维数。

两者之和等于输入空间的维数。
\end{intuition}

% ============================================
\section{非方阵}
% ============================================

\begin{definition}[不同维度之间的变换]
$m \times n$ 矩阵表示从 $\R^n$ 到 $\R^m$ 的变换:
\begin{itemize}
    \item $m > n$: 从低维到高维的变换
    \item $m < n$: 从高维到低维的变换
\end{itemize}
\end{definition}

\begin{example}
$2 \times 3$ 矩阵:
\[
\mA = \twomat{1}{0}{2}{0}{1}{-1}
\]
是从 $\Rthree$ 到 $\Rtwo$ 的变换。它把三维空间``投影''到一个平面上。
\end{example}

% ============================================
\section{方程组解的不同情形}
% ============================================

\begin{theorem}[$\mA\vx = \vb$ 的分析]
\begin{enumerate}
    \item \textbf{唯一解:} 若 $\det(\mA) \neq 0$
    \item \textbf{无穷多解:} 若 $\det(\mA) = 0$ 且 $\vb \in \Col(\mA)$
    \item \textbf{无解:} 若 $\vb \notin \Col(\mA)$
\end{enumerate}
\end{theorem}

\begin{center}
\begin{tikzpicture}[scale=0.9]
    % 情形一: 唯一解
    \begin{scope}[shift={(0,0)}]
        \fill[blue!10] (-1.5,-1.5) rectangle (1.5,1.5);
        \draw[axis] (-1.5,0) -- (1.5,0);
        \draw[axis] (0,-1.5) -- (0,1.5);
        \fill[red] (0.8,0.6) circle (2pt) node[above right] {$\vx$};
        \node at (0,-2) {唯一解};
        \node at (0,-2.5) {\small $\det \neq 0$};
    \end{scope}

    % 情形二: 无穷多解
    \begin{scope}[shift={(5,0)}]
        \draw[blue!30, very thick] (-1.5,-0.75) -- (1.5,0.75);
        \draw[axis] (-1.5,0) -- (1.5,0);
        \draw[axis] (0,-1.5) -- (0,1.5);
        \node at (0,-2) {无穷多解};
        \node at (0,-2.5) {\small $\det = 0$, $\vb \in \Col$};
    \end{scope}

    % 情形三: 无解
    \begin{scope}[shift={(10,0)}]
        \draw[blue!30, very thick] (-1.5,-0.75) -- (1.5,0.75);
        \draw[axis] (-1.5,0) -- (1.5,0);
        \draw[axis] (0,-1.5) -- (0,1.5);
        \fill[red] (0.5,1.2) circle (2pt) node[right] {$\vb$};
        \node[red] at (0.5,0.8) {$\times$};
        \node at (0,-2) {无解};
        \node at (0,-2.5) {\small $\vb \notin \Col$};
    \end{scope}
\end{tikzpicture}
\end{center}

% ============================================
\section{练习}
% ============================================

\begin{exercise}
计算以下矩阵的逆:
\[
\mA = \twomat{2}{5}{1}{3}
\]
\end{exercise}

\begin{exercise}
用逆矩阵解以下方程组:
\[
\system{
2x + 3y &= 7 \\
x + 2y &= 4
}
\]
\end{exercise}

\begin{exercise}
求以下矩阵的零空间:
\[
\mA = \twomat{1}{-2}{-3}{6}
\]
\end{exercise}

\begin{exercise}
求以下矩阵的秩:
\[
\mA = \threemat{1}{2}{3}{2}{4}{6}{1}{1}{1}
\]
\end{exercise}

\begin{exercise}[挑战]
证明 $(\mA\mB)\inv = \mB\inv\mA\inv$ (假设逆都存在)。
\end{exercise}

\begin{problem}
$k$ 取何值时以下矩阵不可逆?
\[
\mA = \twomat{k}{2}{3}{k}
\]
\end{problem}


% 第 6 章: 点积与对偶性
% lecture06.tex - 点积与对偶性
% 第 6 章: 点积与对偶性 (Dot Products and Duality)

\chapter{点积与对偶性}
\label{ch:dotproduct}

\begin{abstract}
点积 (dot product) 是线性代数中最基本的运算之一。在本章中, 我们将学习它的定义、性质, 以及把向量与线性函数联系起来的深刻概念---对偶性。
\end{abstract}

% ============================================
\section{点积的定义}
% ============================================

\begin{definition}[点积]
两个向量 $\vv = \twovec{v_1}{v_2}$ 和 $\vw = \twovec{w_1}{w_2}$ 的\vocab{点积 (dot product)}为:
\[
\vv \cdot \vw = v_1 w_1 + v_2 w_2
\]
在 $n$ 维空间中:
\[
\vv \cdot \vw = \sum_{i=1}^{n} v_i w_i
\]
\end{definition}

\begin{important}
两个向量的点积是一个\textbf{数} (标量), 不是向量!
\end{important}

\begin{example}
\[
\twovec{3}{2} \cdot \twovec{1}{4} = 3 \times 1 + 2 \times 4 = 3 + 8 = 11
\]
\end{example}

% ============================================
\section{几何解释}
% ============================================

\begin{theorem}[点积的几何公式]
\[
\vv \cdot \vw = \norm{\vv} \norm{\vw} \cos\theta
\]
其中 $\theta$ 是两向量之间的夹角。
\end{theorem}

\begin{intuition}
点积有两种理解方式:
\begin{enumerate}
    \item $\vv$ 的长度乘以 $\vw$ 在 $\vv$ 上的投影 (projection) 的长度
    \item $\vw$ 的长度乘以 $\vv$ 在 $\vw$ 上的投影的长度
\end{enumerate}
\end{intuition}

\begin{center}
\begin{tikzpicture}[scale=1.5]
    % 向量
    \draw[vector, blue!70, very thick] (0,0) -- (3,0) node[right] {$\vv$};
    \draw[vector, green!60!black, very thick] (0,0) -- (2,1.5) node[above] {$\vw$};

    % 投影
    \draw[dashed, red] (2,1.5) -- (2,0);
    \draw[red, thick] (0,0) -- (2,0) node[midway, below] {$\vw$ 的投影};

    % 夹角
    \draw (0.7,0) arc (0:36.87:0.7) node[midway, right] {$\theta$};
\end{tikzpicture}
\end{center}

% ============================================
\section{点积的性质}
% ============================================

\begin{theorem}[主要性质]
\begin{enumerate}
    \item \textbf{交换律:} $\vv \cdot \vw = \vw \cdot \vv$
    \item \textbf{分配律:} $\vv \cdot (\vw + \vu) = \vv \cdot \vw + \vv \cdot \vu$
    \item \textbf{数乘:} $(c\vv) \cdot \vw = c(\vv \cdot \vw)$
    \item \textbf{正定性:} $\vv \cdot \vv \geq 0$ 且 $\vv \cdot \vv = 0 \Leftrightarrow \vv = \vzero$
\end{enumerate}
\end{theorem}

\begin{definition}[向量的长度]
\[
\norm{\vv} = \sqrt{\vv \cdot \vv} = \sqrt{v_1^2 + v_2^2 + \cdots + v_n^2}
\]
\end{definition}

% ============================================
\section{正交性}
% ============================================

\begin{theorem}[正交的条件]
两个向量 $\vv$ 和 $\vw$ \vocab{正交 (orthogonal)} (垂直) 当且仅当:
\[
\vv \cdot \vw = 0
\]
\end{theorem}

\begin{intuition}
若 $\vv \cdot \vw = 0$, 则 $\cos\theta = 0$, 所以 $\theta = 90°$。

任何向量在与它垂直的向量上的投影都是零。
\end{intuition}

\begin{example}
向量 $\vv = \twovec{3}{2}$ 和 $\vw = \twovec{2}{-3}$ 正交, 因为:
\[
\vv \cdot \vw = 3 \times 2 + 2 \times (-3) = 6 - 6 = 0
\]
\end{example}

% ============================================
\section{向量投影}
% ============================================

\begin{definition}[向量的投影]
向量 $\vv$ 在向量 $\vw$ 上的投影为:
\[
\proj_{\vw}(\vv) = \frac{\vv \cdot \vw}{\vw \cdot \vw} \vw = \frac{\vv \cdot \vw}{\norm{\vw}^2} \vw
\]
\end{definition}

\begin{intuition}
$\vv$ 在 $\vw$ 上的投影是一个沿 $\vw$ 方向的向量, 它表示 $\vv$ 在 $\vw$ 所在直线上的``影子''。
\end{intuition}

\begin{center}
\begin{tikzpicture}[scale=1.5]
    \draw[vector, blue!70, very thick] (0,0) -- (3,0) node[right] {$\vw$};
    \draw[vector, green!60!black, very thick] (0,0) -- (2,2) node[above] {$\vv$};
    \draw[vector, red!70, very thick] (0,0) -- (2,0) node[below] {$\proj_{\vw}(\vv)$};
    \draw[dashed, gray] (2,2) -- (2,0);
\end{tikzpicture}
\end{center}

% ============================================
\section{对偶性}
% ============================================

\begin{intuition}
一个深刻的洞见: 每个 $1 \times n$ 行向量都可以看作一个\textbf{线性函数}, 它把 $n$ 维向量变成数。

例如, $\bmat{2 & 1}$ 是一个线性函数:
\[
\bmat{2 & 1} \twovec{x}{y} = 2x + y
\]
\end{intuition}

\begin{theorem}[对偶性]
每个线性函数 $f: \R^n \to \R$ 都可以写成与某个固定向量的点积:
\[
f(\vx) = \vv \cdot \vx
\]
其中向量 $\vv$ 唯一确定。
\end{theorem}

\begin{definition}[对偶向量]
表示线性函数 $f$ 的向量 $\vv$ 称为该函数的\vocab{对偶向量 (dual vector)}。
\end{definition}

\begin{practical}
\textbf{在机器学习中的应用:}

在神经网络中, 每个神经元计算其输入的一个线性函数:
\[
\text{output} = w_1 x_1 + w_2 x_2 + \cdots + w_n x_n = \vw \cdot \vx
\]
权重 $\vw$ 就是该神经元的对偶向量。
\end{practical}

% ============================================
\section{实际应用}
% ============================================

\begin{practical}
\textbf{物理: 机械功}

力 $\vec{F}$ 在位移 $\vec{d}$ 上做的功为:
\[
W = \vec{F} \cdot \vec{d} = \norm{\vec{F}} \norm{\vec{d}} \cos\theta
\]
若力的方向与运动方向垂直, 则不做功!
\end{practical}

\begin{practical}
\textbf{计算机图形学: 光照计算}

表面反射的光强为:
\[
I = \max(0, \vec{n} \cdot \vec{l})
\]
其中 $\vec{n}$ 是表面法向量, $\vec{l}$ 是光照方向。
\end{practical}

\begin{practical}
\textbf{文本相似度}

要比较两篇文本文档, 把每篇转换为向量 (例如 TF-IDF), 再计算余弦相似度:
\[
\text{similarity} = \frac{\vv \cdot \vw}{\norm{\vv}\norm{\vw}} = \cos\theta
\]
\end{practical}

% ============================================
\section{练习}
% ============================================

\begin{exercise}
计算以下点积:
\begin{enumerate}[label=(\alph*)]
    \item $\twovec{1}{2} \cdot \twovec{3}{4}$
    \item $\threevec{1}{-1}{2} \cdot \threevec{2}{3}{-1}$
\end{enumerate}
\end{exercise}

\begin{exercise}
向量 $\va = \twovec{4}{3}$ 和 $\vb = \twovec{-3}{4}$ 正交吗?
\end{exercise}

\begin{exercise}
求 $\vv = \twovec{3}{4}$ 在 $\vw = \twovec{1}{0}$ 上的投影。
\end{exercise}

\begin{exercise}
计算向量 $\va = \twovec{1}{1}$ 与 $\vb = \twovec{1}{0}$ 之间的夹角。
\end{exercise}

\begin{exercise}[挑战]
证明对任何两个向量 $\vv$ 和 $\vw$:
\[
\norm{\vv + \vw}^2 = \norm{\vv}^2 + 2(\vv \cdot \vw) + \norm{\vw}^2
\]
\end{exercise}

\begin{problem}
一个大小为 $10$ 牛顿的力以与水平方向成 $60°$ 夹角作用在物体上。若物体水平移动了 $5$ 米, 做了多少功?
\end{problem}


% 第 7 章: 叉积及应用
% lecture07.tex - 叉积及应用
% 第 7 章: 叉积及应用 (Cross Products and Applications)

\chapter{叉积及应用}
\label{ch:crossproduct}

\begin{abstract}
叉积 (cross product, 又称向量积 vector product) 是一种运算, 它由三维空间中的两个向量生成一个同时垂直于这两个向量的向量。在本章中, 我们将学习它的定义、性质, 以及它与行列式之间的深刻联系。
\end{abstract}

% ============================================
\section{二维空间中的叉积}
% ============================================

\begin{definition}[二维叉积]
对于 $\Rtwo$ 中的两个向量:
\[
\vv \times \vw = v_1 w_2 - v_2 w_1 = \det\twomat{v_1}{w_1}{v_2}{w_2}
\]
结果是一个\textbf{数 (number)}, 而不是向量。
\end{definition}

\begin{intuition}
二维叉积 (2D cross product) = 由两个向量张成的平行四边形的有向面积

正号: $\vw$ 在 $\vv$ 的左侧\\
负号: $\vw$ 在 $\vv$ 的右侧
\end{intuition}

\begin{center}
\begin{tikzpicture}[scale=1.2]
    % 平行四边形
    \fill[blue!20] (0,0) -- (2,0.5) -- (3,2) -- (1,1.5) -- cycle;
    \draw[vector, red!70, very thick] (0,0) -- (2,0.5) node[midway, below] {$\vv$};
    \draw[vector, green!60!black, very thick] (0,0) -- (1,1.5) node[midway, left] {$\vw$};
    \draw[dashed] (2,0.5) -- (3,2);
    \draw[dashed] (1,1.5) -- (3,2);
    \node at (1.5,1) {面积 $= |\vv \times \vw|$};
\end{tikzpicture}
\end{center}

% ============================================
\section{三维空间中的叉积}
% ============================================

\begin{definition}[三维叉积]
对于 $\vv = \threevec{v_1}{v_2}{v_3}$ 和 $\vw = \threevec{w_1}{w_2}{w_3}$:
\[
\vv \times \vw = \threevec{v_2 w_3 - v_3 w_2}{v_3 w_1 - v_1 w_3}{v_1 w_2 - v_2 w_1}
\]
\end{definition}

\begin{theorem}[行列式公式]
\[
\vv \times \vw = \det\threemat{\vi}{\vj}{\vk}{v_1}{v_2}{v_3}{w_1}{w_2}{w_3}
\]
(按第一行作形式展开)
\end{theorem}

\begin{intuition}
叉积 $\vv \times \vw$:
\begin{itemize}
    \item \textbf{方向 (direction):} 同时垂直于两个向量 (右手定则 right-hand rule)
    \item \textbf{大小 (magnitude):} 由 $\vv$ 和 $\vw$ 张成的平行四边形的面积
\end{itemize}
\end{intuition}

\begin{center}
\begin{tikzpicture}[scale=1]
    % 三维坐标示意
    \draw[axis] (0,0) -- (3,0) node[right] {$x$};
    \draw[axis] (0,0) -- (0,3) node[above] {$z$};
    \draw[axis] (0,0) -- (-1,-0.7) node[below left] {$y$};

    % 向量
    \draw[vector, red!70, very thick] (0,0) -- (2,-0.5) node[right] {$\vv$};
    \draw[vector, green!60!black, very thick] (0,0) -- (-0.5,0) node[left] {$\vw$};
    \draw[vector, blue!70, ultra thick] (0,0) -- (0,2) node[above] {$\vv \times \vw$};

    % 平行四边形底面
    \fill[gray!20] (0,0) -- (2,-0.5) -- (1.5,-0.5) -- (-0.5,0) -- cycle;
\end{tikzpicture}
\end{center}

% ============================================
\section{叉积的性质}
% ============================================

\begin{theorem}[主要性质]
\begin{enumerate}
    \item \textbf{反交换律 (anti-commutativity):} $\vv \times \vw = -(\vw \times \vv)$
    \item \textbf{分配律 (distributivity):} $\vv \times (\vw + \vu) = \vv \times \vw + \vv \times \vu$
    \item \textbf{标量乘法 (scalar multiplication):} $(c\vv) \times \vw = c(\vv \times \vw)$
    \item \textbf{正交性 (orthogonality):} $\vv \times \vw \perp \vv$ 且 $\vv \times \vw \perp \vw$
    \item \textbf{自身叉积为零 (self-cross is zero):} $\vv \times \vv = \vzero$
\end{enumerate}
\end{theorem}

\begin{warning}
叉积\textbf{不满足交换律}!
\[
\vv \times \vw \neq \vw \times \vv
\]
事实上, 它们互为相反向量。
\end{warning}

\begin{theorem}[模长公式]
\[
\norm{\vv \times \vw} = \norm{\vv} \norm{\vw} \sin\theta
\]
其中 $\theta$ 是两个向量之间的夹角。
\end{theorem}

% ============================================
\section{从线性变换视角看叉积}
% ============================================

\begin{intuition}
更深刻的看法: 叉积 (cross product) 可以通过对偶 (duality) 来定义。

函数 $f(\vx) = \det[\vv \,|\, \vw \,|\, \vx]$ 是关于 $\vx$ 的线性函数 (linear function)。由对偶性, 必存在一个向量 $\vp$ 使得:
\[
f(\vx) = \vp \cdot \vx
\]
这个 $\vp$ 正是 $\vv \times \vw$!
\end{intuition}

\begin{theorem}
\[
\det[\vv \,|\, \vw \,|\, \vx] = (\vv \times \vw) \cdot \vx
\]
\end{theorem}

% ============================================
\section{基向量的叉积}
% ============================================

\begin{theorem}[基向量的叉积]
\begin{align*}
\vi \times \vj &= \vk & \vj \times \vi &= -\vk \\
\vj \times \vk &= \vi & \vk \times \vj &= -\vi \\
\vk \times \vi &= \vj & \vi \times \vk &= -\vj
\end{align*}
\end{theorem}

\begin{example}
\[
\threevec{2}{3}{4} \times \threevec{5}{6}{7}
\]
\begin{align*}
&= (3 \times 7 - 4 \times 6)\vi - (2 \times 7 - 4 \times 5)\vj + (2 \times 6 - 3 \times 5)\vk \\
&= (21-24)\vi - (14-20)\vj + (12-15)\vk \\
&= -3\vi + 6\vj - 3\vk = \threevec{-3}{6}{-3}
\end{align*}
\end{example}

% ============================================
\section{应用}
% ============================================

\begin{practical}
\textbf{物理: 力矩 (torque)}

力 $\vec{F}$ 关于距离为 $\vec{r}$ 的一点的力矩为:
\[
\vec{\tau} = \vec{r} \times \vec{F}
\]
力矩的方向垂直于力与力臂所在的平面。
\end{practical}

\begin{practical}
\textbf{物理: 洛伦兹力 (Lorentz force)}

在磁场中运动的带电粒子所受的力为:
\[
\vec{F} = q\vec{v} \times \vec{B}
\]
\end{practical}

\begin{practical}
\textbf{计算机图形学: 曲面法向量 (surface normal)}

求顶点为 $A$, $B$, $C$ 的三角形面的法向量 (normal vector):
\[
\vec{n} = (\vec{B} - \vec{A}) \times (\vec{C} - \vec{A})
\]
\end{practical}

\begin{practical}
\textbf{三角形面积计算}

顶点为 $A$, $B$, $C$ 的三角形的面积为:
\[
\text{Area} = \frac{1}{2}\norm{(\vec{B}-\vec{A}) \times (\vec{C}-\vec{A})}
\]
\end{practical}

% ============================================
\section{三重积}
% ============================================

\begin{definition}[标量三重积 (scalar triple product)]
\[
\va \cdot (\vb \times \vc) = \det[\va \,|\, \vb \,|\, \vc]
\]
结果 = 由三个向量张成的平行六面体的有向体积
\end{definition}

\begin{theorem}[轮换性质]
\[
\va \cdot (\vb \times \vc) = \vb \cdot (\vc \times \va) = \vc \cdot (\va \times \vb)
\]
\end{theorem}

% ============================================
\section{练习}
% ============================================

\begin{exercise}
计算叉积:
\[
\va = \threevec{1}{2}{3}, \quad \vb = \threevec{4}{5}{6}
\]
\end{exercise}

\begin{exercise}
求以 $\va = \threevec{1}{0}{0}$ 和 $\vb = \threevec{1}{1}{0}$ 为边的平行四边形的面积。
\end{exercise}

\begin{exercise}
证明 $\va \times \vb$ 同时垂直于 $\va$ 和 $\vb$。
\end{exercise}

\begin{exercise}
求过点 $(1,0,0)$, $(0,1,0)$, $(0,0,1)$ 的平面的法向量。
\end{exercise}

\begin{exercise}[挑战]
证明:
\[
\norm{\va \times \vb}^2 + (\va \cdot \vb)^2 = \norm{\va}^2 \norm{\vb}^2
\]
(这称为拉格朗日恒等式 Lagrange's identity)
\end{exercise}

\begin{problem}
计算以以下向量为棱的平行六面体的体积:
\[
\va = \threevec{1}{1}{0}, \quad \vb = \threevec{0}{1}{1}, \quad \vc = \threevec{1}{0}{1}
\]
\end{problem}


% 第 8 章: 克拉默法则
% lecture08.tex - 克拉默法则
% 第 8 章: 克拉默法则 (Cramer's Rule)

\chapter{克拉默法则}
\label{ch:cramer}

\begin{abstract}
克拉默法则 (Cramer's Rule) 是一种用行列式求解线性方程组的优雅方法。在本章中, 我们将学习它的几何解释及其使用条件。
\end{abstract}

% ============================================
\section{克拉默法则的表述}
% ============================================

\begin{theorem}[克拉默法则]
对于系数矩阵为 $n \times n$ 矩阵且 $\det(\mA) \neq 0$ 的方程组 $\mA\vx = \vb$:
\[
x_i = \frac{\det(\mA_i)}{\det(\mA)}
\]
其中 $\mA_i$ 是把 $\mA$ 的第 $i$ 列换成向量 $\vb$ 后得到的矩阵。
\end{theorem}

\begin{example}
对于 $2 \times 2$ 方程组:
\[
\system{
ax + by &= e \\
cx + dy &= f
}
\]
解为:
\[
x = \frac{\det\twomat{e}{b}{f}{d}}{\det\twomat{a}{b}{c}{d}} = \frac{ed - bf}{ad - bc}
\]
\[
y = \frac{\det\twomat{a}{e}{c}{f}}{\det\twomat{a}{b}{c}{d}} = \frac{af - ec}{ad - bc}
\]
\end{example}

% ============================================
\section{几何解释}
% ============================================

\begin{intuition}
在二维情形, 方程 $\mA\vx = \vb$ 是在问:

``$\mA$ 的列向量的什么线性组合 (linear combination) 等于 $\vb$?''

若 $\mA = [\va_1 \,|\, \va_2]$ 且 $\vx = \twovec{x}{y}$:
\[
x\va_1 + y\va_2 = \vb
\]
\end{intuition}

\begin{center}
\begin{tikzpicture}[scale=1.3]
    % 平行四边形向量
    \draw[vector, red!70, very thick] (0,0) -- (2,0.5) node[right] {$\va_1$};
    \draw[vector, blue!70, very thick] (0,0) -- (0.5,2) node[above] {$\va_2$};
    \draw[vector, green!60!black, ultra thick] (0,0) -- (1.5,1.5) node[above right] {$\vb$};

    % 由 b 和 a2 张成的平行四边形
    \fill[orange!20] (0,0) -- (1.5,1.5) -- (2,3.5) -- (0.5,2) -- cycle;

    % 由 a1 和 a2 张成的平行四边形
    \fill[gray!20] (0,0) -- (2,0.5) -- (2.5,2.5) -- (0.5,2) -- cycle;

    \node at (1.2,2.5) {\small 橙色区域};
    \node at (1.8,1) {\small 灰色区域};
\end{tikzpicture}
\end{center}

\begin{theorem}[面积解释]
\[
x = \frac{\text{Area of parallelogram}(\vb, \va_2)}{\text{Area of parallelogram}(\va_1, \va_2)}
\]
\[
y = \frac{\text{Area of parallelogram}(\va_1, \vb)}{\text{Area of parallelogram}(\va_1, \va_2)}
\]
\end{theorem}

\begin{intuition}
为什么这样是对的?

考虑平行四边形 $(\vb, \va_2)$。因为 $\vb = x\va_1 + y\va_2$:
\[
\text{Area}(\vb, \va_2) = \text{Area}(x\va_1 + y\va_2, \va_2)
\]
因为 $\va_2 \times \va_2 = 0$:
\[
= \text{Area}(x\va_1, \va_2) = x \cdot \text{Area}(\va_1, \va_2)
\]
因此:
\[
x = \frac{\text{Area}(\vb, \va_2)}{\text{Area}(\va_1, \va_2)}
\]
\end{intuition}

% ============================================
\section{计算示例}
% ============================================

\begin{example}
解方程组:
\[
\system{
3x + 2y &= 7 \\
x + 4y &= 9
}
\]

\textbf{第 1 步:} 计算 $\det(\mA)$:
\[
\det\twomat{3}{2}{1}{4} = 12 - 2 = 10
\]

\textbf{第 2 步:} 计算 $x$:
\[
x = \frac{\det\twomat{7}{2}{9}{4}}{\det(\mA)} = \frac{28 - 18}{10} = \frac{10}{10} = 1
\]

\textbf{第 3 步:} 计算 $y$:
\[
y = \frac{\det\twomat{3}{7}{1}{9}}{\det(\mA)} = \frac{27 - 7}{10} = \frac{20}{10} = 2
\]

\textbf{检验:} $3(1) + 2(2) = 7$ \checkmark\ 且 $1(1) + 4(2) = 9$ \checkmark
\end{example}

% ============================================
\section{三维及更高维情形}
% ============================================

\begin{example}
一个 $3 \times 3$ 方程组:
\[
\system{
2x + y - z &= 3 \\
x - y + 2z &= 1 \\
3x + 2y + z &= 4
}
\]

系数矩阵为:
\[
\mA = \threemat{2}{1}{-1}{1}{-1}{2}{3}{2}{1}
\]

\[
x = \frac{\det\threemat{3}{1}{-1}{1}{-1}{2}{4}{2}{1}}{\det(\mA)}, \quad
y = \frac{\det\threemat{2}{3}{-1}{1}{1}{2}{3}{4}{1}}{\det(\mA)}, \quad
z = \frac{\det\threemat{2}{1}{3}{1}{-1}{1}{3}{2}{4}}{\det(\mA)}
\]
\end{example}

\begin{intuition}
在三维情形, 我们得到的是\textbf{体积 (volumes)}之比, 而不是面积之比。

$x$ = 平行六面体 $(\vb, \va_2, \va_3)$ 的体积与 $(\va_1, \va_2, \va_3)$ 的体积之比
\end{intuition}

% ============================================
\section{局限与应用}
% ============================================

\begin{warning}
克拉默法则 (Cramer's Rule):
\begin{itemize}
    \item 仅适用于方方程组 ($n$ 个方程, $n$ 个未知数)
    \item 仅当 $\det(\mA) \neq 0$ (有唯一解)
    \item 当 $n$ 较大时, 计算量\textbf{非常大}
\end{itemize}
\end{warning}

\begin{remark}
实际中, 求解大型方程组时一般使用高斯消元法 (Gaussian elimination) 等更高效的方法。但克拉默法则 (Cramer's Rule):
\begin{itemize}
    \item 提供了更深刻的理论理解
    \item 便于给出解析公式
    \item 常用于定理证明
\end{itemize}
\end{remark}

\begin{practical}
\textbf{应用: 求直线交点}

两条直线 $a_1x + b_1y = c_1$ 和 $a_2x + b_2y = c_2$ 的交点为:
\[
x = \frac{c_1 b_2 - c_2 b_1}{a_1 b_2 - a_2 b_1}, \quad y = \frac{a_1 c_2 - a_2 c_1}{a_1 b_2 - a_2 b_1}
\]
这个公式直接来自克拉默法则。
\end{practical}

% ============================================
\section{练习}
% ============================================

\begin{exercise}
用克拉默法则求解:
\[
\system{
2x + 3y &= 8 \\
4x - y &= 2
}
\]
\end{exercise}

\begin{exercise}
用克拉默法则求解:
\[
\system{
x + y + z &= 6 \\
2x - y + z &= 3 \\
x + 2y - z &= 2
}
\]
\end{exercise}

\begin{exercise}
用克拉默法则求直线 $2x + 3y = 7$ 和 $x - y = 1$ 的交点。
\end{exercise}

\begin{exercise}
解释当 $\det(\mA) = 0$ 时克拉默法则的几何含义。
\end{exercise}

\begin{exercise}[挑战]
证明克拉默公式与计算 $\mA\inv \vb$ 是一致的。
\end{exercise}

\begin{problem}
已知三角形顶点为 $A(1,1)$, $B(4,2)$, $C(2,5)$。用克拉默法则求其重心的坐标 (coordinates)。
\end{problem}


% 第 9 章: 基变换
% lecture09.tex - 基变换
% 第 9 章: 基变换 (Change of Basis)

\chapter{基变换}
\label{ch:changebasis}

\begin{abstract}
同一个向量在不同基下有不同的坐标 (coordinates)。在本章中, 我们将学习如何在不同坐标系之间转换, 以及这一概念与矩阵的关系。
\end{abstract}

% ============================================
\section{不同基下的坐标}
% ============================================

\begin{intuition}
向量的坐标 (coordinates) 依赖于所选的基。如果你说不同的语言, 同一个意思也会有不同的表达方式。

例如: 在标准基下为 $\twovec{3}{2}$ 的向量 $\vv$, 在另一组基下可能是 $\twovec{1}{1}$!
\end{intuition}

\begin{definition}[基下的坐标]
若 $\mathcal{B} = \{\vb_1, \vb_2\}$ 是一组基, 则向量 $\vv$ 在该基下的坐标 (coordinates) 是满足下式的系数 $c_1, c_2$:
\[
\vv = c_1 \vb_1 + c_2 \vb_2
\]
记号: $[\vv]_{\mathcal{B}} = \twovec{c_1}{c_2}$
\end{definition}

\begin{example}
基 $\mathcal{B} = \left\{ \vb_1 = \twovec{2}{1}, \vb_2 = \twovec{-1}{1} \right\}$

把向量 $\vv = \twovec{3}{2}$ 表示在这组基下。

需要求 $c_1, c_2$ 使得:
\[
c_1 \twovec{2}{1} + c_2 \twovec{-1}{1} = \twovec{3}{2}
\]
\[
\system{
2c_1 - c_2 &= 3 \\
c_1 + c_2 &= 2
}
\]
解为: $c_1 = \frac{5}{3}$, $c_2 = \frac{1}{3}$

于是 $[\vv]_{\mathcal{B}} = \twovec{5/3}{1/3}$
\end{example}

% ============================================
\section{过渡矩阵}
% ============================================

\begin{definition}[过渡矩阵]
从基 $\mathcal{B}$ 到标准基的\vocab{过渡矩阵 (transition matrix, 即基变换矩阵 change of basis matrix)}是以 $\mathcal{B}$ 的基向量为列构成的矩阵:
\[
\mP = [\vb_1 \,|\, \vb_2 \,|\, \cdots \,|\, \vb_n]
\]
\end{definition}

\begin{theorem}
\[
\vv = \mP \, [\vv]_{\mathcal{B}}
\]
即: 基 $\mathcal{B}$ 下的坐标乘以过渡矩阵 = 标准基下的向量
\end{theorem}

\begin{theorem}
\[
[\vv]_{\mathcal{B}} = \mP\inv \vv
\]
即: 要从标准基转换到基 $\mathcal{B}$, 需要用逆矩阵。
\end{theorem}

% ============================================
\section{基变换的变换视角}
% ============================================

\begin{intuition}
矩阵 $\mP$ 是一个恒等变换, 它只改变语言:
\begin{itemize}
    \item $\mP$: 把语言 $\mathcal{B}$ 翻译成标准语言
    \item $\mP\inv$: 把标准语言翻译成语言 $\mathcal{B}$
\end{itemize}
\end{intuition}

\begin{center}
\begin{tikzpicture}[scale=1.2]
    % 标准基
    \begin{scope}[shift={(0,0)}]
        \draw[grid] (-0.5,-0.5) grid (3.5,2.5);
        \draw[axis] (-0.5,0) -- (3.5,0);
        \draw[axis] (0,-0.5) -- (0,2.5);
        \draw[basis] (0,0) -- (1,0) node[below] {$\vi$};
        \draw[basis] (0,0) -- (0,1) node[left] {$\vj$};
        \draw[vector, blue!70, very thick] (0,0) -- (3,2) node[above] {$\vv$};
        \node at (1.5,-1) {标准基};
        \node at (1.5,-1.5) {$\vv = \twovec{3}{2}$};
    \end{scope}

    % 另一组基
    \begin{scope}[shift={(6,0)}]
        \draw[gray!30] (-0.5,-0.5) -- (1,0.5) -- (3,1.5) -- (1.5,1) -- cycle;
        \draw[axis] (-0.5,0) -- (3.5,0);
        \draw[axis] (0,-0.5) -- (0,2.5);
        \draw[red!70, very thick, -{Stealth}] (0,0) -- (2,1) node[below right] {$\vb_1$};
        \draw[red!70, very thick, -{Stealth}] (0,0) -- (-1,1) node[above left] {$\vb_2$};
        \draw[vector, blue!70, very thick] (0,0) -- (3,2) node[above] {$\vv$};
        \node at (1.5,-1) {基 $\mathcal{B}$};
        \node at (1.5,-1.5) {$[\vv]_{\mathcal{B}} = \twovec{?}{?}$};
    \end{scope}
\end{tikzpicture}
\end{center}

% ============================================
\section{在不同基下表示变换}
% ============================================

\begin{theorem}[新基下的变换矩阵]
若 $\mA$ 是变换 $T$ 在标准基下的矩阵, 则同一变换在基 $\mathcal{B}$ 下的矩阵为:
\[
[\mA]_{\mathcal{B}} = \mP\inv \mA \mP
\]
\end{theorem}

\begin{intuition}
这个公式包含三步:
\begin{enumerate}
    \item $\mP$: 从基 $\mathcal{B}$ 翻译到标准基
    \item $\mA$: 在标准基下实施变换
    \item $\mP\inv$: 把结果翻译回基 $\mathcal{B}$
\end{enumerate}
\end{intuition}

\begin{center}
\begin{tikzpicture}[node distance=3cm]
    \node (B1) {$[\vv]_{\mathcal{B}}$};
    \node (S1) [right of=B1] {$\vv$};
    \node (S2) [right of=S1] {$T(\vv)$};
    \node (B2) [right of=S2] {$[T(\vv)]_{\mathcal{B}}$};

    \draw[-{Stealth}, thick] (B1) -- (S1) node[midway, above] {$\mP$};
    \draw[-{Stealth}, thick] (S1) -- (S2) node[midway, above] {$\mA$};
    \draw[-{Stealth}, thick] (S2) -- (B2) node[midway, above] {$\mP\inv$};
    \draw[-{Stealth}, thick, blue] (B1) to[bend right=30] node[midway, below] {$\mP\inv\mA\mP$} (B2);
\end{tikzpicture}
\end{center}

% ============================================
\section{基变换的重要性}
% ============================================

\begin{intuition}
为什么基变换 (change of basis) 重要?

有些变换在特定基下看起来\textbf{更简单}。例如:
\begin{itemize}
    \item 在标准基下旋转很复杂
    \item 但在一条轴落在旋转轴上的基下, 它就变简单了
\end{itemize}

对一个变换最好的基? \textbf{特征基 (eigenbasis)}---下一章见!
\end{intuition}

\begin{practical}
\textbf{应用: 简化计算}

假设要计算 $\mA^{100}$。如果在合适的基下, $\mA$ 变为对角矩阵:
\[
\mP\inv \mA \mP = \mD \quad \text{(diagonal)}
\]
则:
\[
\mA^{100} = \mP \mD^{100} \mP\inv
\]
而 $\mD^{100}$ 非常容易计算!
\end{practical}

% ============================================
\section{练习}
% ============================================

\begin{exercise}
已知基 $\mathcal{B} = \left\{ \twovec{1}{1}, \twovec{1}{-1} \right\}$, 求向量 $\vv = \twovec{3}{1}$ 在该基下的坐标 (coordinates)。
\end{exercise}

\begin{exercise}
写出从 $\mathcal{B} = \left\{ \twovec{2}{1}, \twovec{1}{1} \right\}$ 到标准基的过渡矩阵 (transition matrix)。
\end{exercise}

\begin{exercise}
若变换在标准基下的矩阵为 $\mA = \twomat{2}{1}{0}{3}$, 求该变换在基 $\mathcal{B} = \left\{ \twovec{1}{0}, \twovec{1}{1} \right\}$ 下的矩阵。
\end{exercise}

\begin{exercise}[挑战]
证明 $\det(\mP\inv\mA\mP) = \det(\mA)$。
\end{exercise}

\begin{problem}
已知两组基 $\mathcal{B}_1$ 和 $\mathcal{B}_2$。如何直接计算从 $\mathcal{B}_1$ 到 $\mathcal{B}_2$ 的过渡矩阵 (不经过标准基)?
\end{problem}


% 第 10 章: 特征向量与特征值
% lecture10.tex - 特征向量与特征值
% 第 10 章: 特征向量与特征值

\chapter{特征向量与特征值}
\label{ch:eigen}

\begin{abstract}
特征向量 (eigenvector) 是线性变换只拉伸而不改变其方向的特殊方向。这一概念是线性代数中最重要的思想之一, 在物理学、工程学和数据科学中有着广泛的应用。
\end{abstract}

% ============================================
\section{动机: 特殊的向量}
% ============================================

\begin{intuition}
当你施加一个线性变换时, 大多数向量都会偏离原来的方向。但有些特殊的向量只被拉伸或压缩, 仍然\textbf{保持在同一条直线上}。

这些特殊的向量称为\textbf{特征向量 (eigenvector)}。
\end{intuition}

\begin{center}
\begin{tikzpicture}[scale=1.2]
    % 变换前
    \begin{scope}[shift={(0,0)}]
        \draw[grid] (-2,-2) grid (2,2);
        \draw[axis] (-2,0) -- (2,0);
        \draw[axis] (0,-2) -- (0,2);
        \draw[vector, blue!70, very thick] (0,0) -- (1,0.5) node[right] {$\vv_1$};
        \draw[vector, red!70, very thick] (0,0) -- (0.5,1) node[above] {$\vv_2$};
        \draw[vector, green!60!black, very thick] (0,0) -- (-1,1) node[above left] {特征方向!};
        \node at (0,-2.5) {变换前};
    \end{scope}

    \draw[-{Stealth}, very thick] (2.5,0) -- (4,0) node[midway, above] {$T$};

    % 变换后
    \begin{scope}[shift={(6.5,0)}]
        \draw[blue!20] (-2,-1) -- (0,2) -- (4,2) -- (2,-1) -- cycle;
        \draw[axis] (-2,0) -- (4,0);
        \draw[axis] (0,-2) -- (0,2);
        \draw[vector, blue!70, very thick] (0,0) -- (1.5,1) node[right] {$T(\vv_1)$};
        \draw[vector, red!70, very thick] (0,0) -- (1.5,1.5) node[above] {$T(\vv_2)$};
        \draw[vector, green!60!black, very thick] (0,0) -- (-2,2) node[above left] {$2 \times$ 特征方向!};
        \node at (1,-2.5) {变换后};
    \end{scope}
\end{tikzpicture}
\end{center}

% ============================================
\section{正式定义}
% ============================================

\begin{definition}[特征向量与特征值]
非零向量 $\vv$ 是矩阵 $\mA$ 的\vocab{特征向量 (eigenvector)}, 如果:
\[
\mA\vv = \lambda\vv
\]
对某个标量 $\lambda$ 成立。该标量称为对应的\vocab{特征值 (eigenvalue)}。
\end{definition}

\begin{intuition}
$\mA\vv = \lambda\vv$ 的意思是:
\begin{center}
``变换 $\mA$ 作用在向量 $\vv$ 上, 就相当于一次标量乘法''
\end{center}
向量 $\vv$ 保持在同一条直线上, 只是被 $\lambda$ 缩放了。
\end{intuition}

\begin{example}
对于矩阵 $\mA = \twomat{3}{1}{0}{2}$:

向量 $\vv = \twovec{1}{0}$ 是一个特征向量:
\[
\mA\vv = \twomat{3}{1}{0}{2}\twovec{1}{0} = \twovec{3}{0} = 3\twovec{1}{0} = 3\vv
\]
对应的特征值为: $\lambda = 3$
\end{example}

% ============================================
\section{求特征值}
% ============================================

\begin{theorem}[特征方程]
$\lambda$ 是 $\mA$ 的特征值当且仅当:
\[
\det(\mA - \lambda\mI) = 0
\]
\end{theorem}

\begin{proof}
$\mA\vv = \lambda\vv$ 可写为:
\[
\mA\vv - \lambda\vv = \vzero \implies (\mA - \lambda\mI)\vv = \vzero
\]
该方程有非零解当且仅当 $\mA - \lambda\mI$ 不可逆, 即 $\det(\mA - \lambda\mI) = 0$。
\end{proof}

\begin{definition}[特征多项式]
$\det(\mA - \lambda\mI)$ 是关于 $\lambda$ 的多项式, 称为\vocab{特征多项式 (characteristic polynomial)}。它的根就是特征值。
\end{definition}

\begin{example}
对于 $\mA = \twomat{3}{1}{0}{2}$:
\[
\det(\mA - \lambda\mI) = \det\twomat{3-\lambda}{1}{0}{2-\lambda} = (3-\lambda)(2-\lambda) - 0 = 0
\]
根为: $\lambda_1 = 3$, $\lambda_2 = 2$
\end{example}

% ============================================
\section{求特征向量}
% ============================================

\begin{theorem}
对于每个特征值 $\lambda$, 对应的特征向量可通过求解齐次方程组得到:
\[
(\mA - \lambda\mI)\vv = \vzero
\]
\end{theorem}

\begin{example}
继续上一例, 取 $\lambda = 2$:
\[
(\mA - 2\mI)\vv = \twomat{1}{1}{0}{0}\twovec{v_1}{v_2} = \vzero
\]
方程为: $v_1 + v_2 = 0$, 于是 $v_1 = -v_2$

特征向量为: $\vv = t\twovec{-1}{1}$, 其中 $t \neq 0$ 为任意数
\end{example}

% ============================================
\section{特征空间}
% ============================================

\begin{definition}[特征空间]
对应于特征值 $\lambda$ 的\vocab{特征空间 (eigenspace)}:
\[
E_\lambda = \Null(\mA - \lambda\mI) = \{\vv \mid \mA\vv = \lambda\vv\}
\]
\end{definition}

\begin{intuition}
特征空间包含变换 $\mA$ 仅以倍数 $\lambda$ 缩放的所有向量。该空间总是一个向量\vocab{子空间 (subspace)}。
\end{intuition}

% ============================================
\section{几何解释}
% ============================================

\begin{practical}
\textbf{三维旋转}

对于 $\Rthree$ 中的旋转, 特征值 $\lambda = 1$ 对应的特征向量就是\textbf{旋转轴}! 该向量保持不动。

用轴和角描述旋转, 比用 $3 \times 3$ 矩阵简洁得多。
\end{practical}

\begin{example}
错切矩阵: $\mA = \twomat{1}{1}{0}{1}$

特征值: $\lambda = 1$ (重根)

特征向量: $\twovec{1}{0}$ (只有一个特征方向)

解释: 错切保持 $x$ 轴不动。
\end{example}

% ============================================
\section{特殊情形}
% ============================================

\begin{theorem}[二维旋转]
旋转矩阵 $\mR_\theta = \twomat{\cos\theta}{-\sin\theta}{\sin\theta}{\cos\theta}$, 当 $\theta \neq 0, \pi$ 时:

特征值为: $\lambda = \cos\theta \pm i\sin\theta$ (复数!)

没有实向量保持不动——所有向量都旋转了。
\end{theorem}

\begin{warning}
即使对实矩阵, 特征值也可能是\textbf{复数}! 旋转中就会出现这种情况。
\end{warning}

% ============================================
\section{应用}
% ============================================

\begin{practical}
\textbf{谷歌 PageRank}

把网页看作向量, 把链接看作矩阵。主特征向量 (对应最大特征值的特征向量) 反映了各网页的相对重要性。
\end{practical}

\begin{practical}
\textbf{主成分分析 (PCA)}

在机器学习中, 协方差矩阵的特征向量给出了数据变化的主要方向。
\end{practical}

\begin{practical}
\textbf{量子力学}

算子的特征值 = 可能的测量结果

特征向量 = 系统的稳定状态
\end{practical}

% ============================================
\section{练习}
% ============================================

\begin{exercise}
求下列矩阵的特征值与特征向量:
\[
\mA = \twomat{4}{1}{2}{3}
\]
\end{exercise}

\begin{exercise}
证明对角矩阵 (diagonal matrix) 的特征值就是它的对角元。
\end{exercise}

\begin{exercise}
求旋转 $90°$ 的旋转矩阵的特征值。
\end{exercise}

\begin{exercise}
若 $\lambda$ 是 $\mA$ 的特征值, 证明 $\lambda^2$ 是 $\mA^2$ 的特征值。
\end{exercise}

\begin{exercise}[挑战]
证明矩阵的迹 (对角元之和) 等于特征值之和。
\end{exercise}

\begin{problem}
种群矩阵:
\[
\mL = \twomat{0}{4}{0.5}{0}
\]
求主特征值并解释其含义。
\end{problem}


% 第 11 章: 特征值计算与特征基
% lecture11.tex - 特征值计算与特征基
% 第 11 章: 特征值计算与特征基

\chapter{特征值计算与特征基}
\label{ch:eigenbasis}

\begin{abstract}
在本章中, 我们将学习计算 $2 \times 2$ 矩阵特征值的快捷方法, 并熟悉特征基 (eigenbasis) 与对角化的概念。
\end{abstract}

% ============================================
\section{$2 \times 2$ 矩阵的快捷算法}
% ============================================

\begin{theorem}[快捷公式]
对于矩阵 $\mA = \twomat{a}{b}{c}{d}$:

\textbf{特征值之和:}
\[
\lambda_1 + \lambda_2 = a + d = \tr(\mA)
\]

\textbf{特征值之积:}
\[
\lambda_1 \cdot \lambda_2 = ad - bc = \det(\mA)
\]
\end{theorem}

\begin{intuition}
已知 $m = \lambda_1 + \lambda_2$ 与 $p = \lambda_1 \lambda_2$, 可求出 $\lambda_{1,2}$:
\[
\lambda = \frac{m}{2} \pm \sqrt{\left(\frac{m}{2}\right)^2 - p}
\]
或者: 找出和为 $m$、积为 $p$ 的两个数。
\end{intuition}

\begin{example}
$\mA = \twomat{3}{1}{4}{1}$

迹: $m = 3 + 1 = 4$

行列式: $p = 3 - 4 = -1$

特征值: 满足 $x + y = 4$ 且 $xy = -1$ 的两个数:
\[
\lambda = 2 \pm \sqrt{4 - (-1)} = 2 \pm \sqrt{5}
\]
\end{example}

% ============================================
\section{特征基}
% ============================================

\begin{definition}[特征基]
若矩阵的特征向量可以构成一组\vocab{基 (basis)} (即有 $n$ 个线性无关 (linear independence) 的特征向量), 则该基称为\vocab{特征基 (eigenbasis)}。
\end{definition}

\begin{theorem}
在特征基下, 变换矩阵变为\vocab{对角的 (diagonal)}:
\[
[\mA]_{\text{eigenbasis}} = \mD = \threemat{\lambda_1}{0}{0}{0}{\lambda_2}{0}{0}{0}{\ddots}
\]
\end{theorem}

\begin{intuition}
为什么是对角的?

在特征基下, 每个基向量只被缩放:
\begin{align*}
\mA \vv_1 &= \lambda_1 \vv_1 \to \text{第一列: } \threevec{\lambda_1}{0}{0} \\
\mA \vv_2 &= \lambda_2 \vv_2 \to \text{第二列: } \threevec{0}{\lambda_2}{0}
\end{align*}
\end{intuition}

% ============================================
\section{对角化}
% ============================================

\begin{definition}[对角化]
若矩阵 $\mA$ 可写为:
\[
\mA = \mP \mD \mP\inv
\]
其中 $\mD$ 为对角矩阵 (diagonal matrix), $\mP$ 为以特征向量为列的矩阵, 则称 $\mA$ \vocab{可对角化 (diagonalizable)}。
\end{definition}

\begin{theorem}[可对角化条件]
$n \times n$ 矩阵可对角化当且仅当它有 $n$ 个线性无关的特征向量。
\end{theorem}

\begin{example}
$\mA = \twomat{3}{1}{0}{2}$, 特征值为 $\lambda_1 = 3$, $\lambda_2 = 2$

特征向量: $\vv_1 = \twovec{1}{0}$, $\vv_2 = \twovec{-1}{1}$

\[
\mP = \twomat{1}{-1}{0}{1}, \quad \mD = \twomat{3}{0}{0}{2}
\]

验证: $\mP\mD\mP\inv = \mA$ \checkmark
\end{example}

% ============================================
\section{用对角化计算矩阵幂}
% ============================================

\begin{theorem}
若 $\mA = \mP\mD\mP\inv$:
\[
\mA^n = \mP \mD^n \mP\inv
\]
而 $\mD^n$ 非常简单:
\[
\mD^n = \threemat{\lambda_1^n}{0}{0}{0}{\lambda_2^n}{0}{0}{0}{\ddots}
\]
\end{theorem}

\begin{example}
计算上一例中 $\mA^{100}$:
\[
\mA^{100} = \mP \twomat{3^{100}}{0}{0}{2^{100}} \mP\inv
\]
若不用对角化, 需要做 100 次矩阵乘法!
\end{example}

\begin{practical}
\textbf{应用: 斐波那契数列}

斐波那契数列: $F_{n+1} = F_n + F_{n-1}$

矩阵形式:
\[
\twomat{1}{1}{1}{0}^n = \twomat{F_{n+1}}{F_n}{F_n}{F_{n-1}}
\]

利用对角化, 可得到闭式公式:
\[
F_n = \frac{1}{\sqrt{5}}\left[\left(\frac{1+\sqrt{5}}{2}\right)^n - \left(\frac{1-\sqrt{5}}{2}\right)^n\right]
\]
\end{practical}

% ============================================
\section{不可对角化的矩阵}
% ============================================

\begin{example}
错切矩阵 $\mA = \twomat{1}{1}{0}{1}$:

特征多项式: $(1-\lambda)^2 = 0$

特征值: $\lambda = 1$ (重根)

特征向量: 仅 $\twovec{1}{0}$ (一维)

该矩阵不可对角化!
\end{example}

\begin{warning}
重特征值本身不一定是问题。当独立特征向量的个数少于该特征值的重数时, 问题才会出现。
\end{warning}

% ============================================
\section{对称矩阵}
% ============================================

\begin{theorem}[谱定理]
对称矩阵 ($\mA = \mA\trans$):
\begin{enumerate}
    \item 有实特征值
    \item 属于不同特征值的特征向量相互正交
    \item 总可对角化 (且有一组正交基)
\end{enumerate}
\end{theorem}

\begin{intuition}
对称矩阵是``表现良好''的矩阵。它们总能化为对角形, 且特征向量相互垂直——就像椭圆的主轴一样。
\end{intuition}

% ============================================
\section{练习}
% ============================================

\begin{exercise}
用快捷公式求下列矩阵的特征值:
\[
\mA = \twomat{5}{2}{2}{2}
\]
\end{exercise}

\begin{exercise}
将矩阵 $\mA = \twomat{2}{1}{1}{2}$ 对角化。
\end{exercise}

\begin{exercise}
利用对角化计算 $\mA^{10}$:
\[
\mA = \twomat{1}{1}{0}{2}
\]
\end{exercise}

\begin{exercise}
下列矩阵是否可对角化?
\[
\mA = \twomat{2}{1}{0}{2}
\]
\end{exercise}

\begin{exercise}[挑战]
证明 $\mA$ 与 $\mA\trans$ 有相同的特征值。
\end{exercise}

\begin{problem}
兔子与狐狸的数量由下式建模:
\[
\twovec{R_{n+1}}{F_{n+1}} = \twomat{1.1}{-0.4}{0.2}{0.8}\twovec{R_n}{F_n}
\]
分析种群的长期行为。
\end{problem}


% 第 12 章: 抽象向量空间
% lecture12.tex - 抽象向量空间
% 第 12 章: 抽象向量空间

\chapter{抽象向量空间}
\label{ch:abstract}

\begin{abstract}
在最后一章中, 我们把向量的概念推广到几何箭头之外。我们将看到, 函数、多项式乃至音乐都可以是``向量''——只要它们遵循线性代数的规则。
\end{abstract}

% ============================================
\section{动机: 什么是向量?}
% ============================================

\begin{intuition}
到目前为止, 我们把向量看作箭头或数表。但数学家追问了一个更深刻的问题:

\textbf{``什么样的东西表现得像向量?''}

答案: 凡是能相加、能与标量相乘的东西!
\end{intuition}

\begin{example}
函数可以相加: $(f + g)(x) = f(x) + g(x)$

函数可以与数相乘: $(cf)(x) = c \cdot f(x)$

所以函数也可以是``向量''!
\end{example}

% ============================================
\section{向量空间的正式定义}
% ============================================

\begin{definition}[向量空间]
实数域 $\R$ 上的\vocab{向量空间 (vector space)}是一个集合 $V$, 带有两种运算:
\begin{itemize}
    \item \textbf{加法:} $+ : V \times V \to V$
    \item \textbf{标量乘法:} $\cdot : \R \times V \to V$
\end{itemize}
且满足以下公理。
\end{definition}

\begin{theorem}[向量空间公理]
对一切 $\vu, \vv, \vw \in V$ 与 $a, b \in \R$:

\textbf{加法公理:}
\begin{enumerate}
    \item $\vu + \vv = \vv + \vu$ (交换律)
    \item $(\vu + \vv) + \vw = \vu + (\vv + \vw)$ (结合律)
    \item 单位元存在: $\exists \vzero : \vv + \vzero = \vv$
    \item 逆元存在: $\forall \vv, \exists (-\vv) : \vv + (-\vv) = \vzero$
\end{enumerate}

\textbf{标量乘法公理:}
\begin{enumerate}[resume]
    \item $a(b\vv) = (ab)\vv$
    \item $1 \cdot \vv = \vv$
\end{enumerate}

\textbf{分配律公理:}
\begin{enumerate}[resume]
    \item $a(\vu + \vv) = a\vu + a\vv$
    \item $(a + b)\vv = a\vv + b\vv$
\end{enumerate}
\end{theorem}

% ============================================
\section{向量空间的例子}
% ============================================

\subsection{$\R^n$——标准空间}

\begin{example}
$\R^n = \{(x_1, x_2, \ldots, x_n) \mid x_i \in \R\}$, 按通常的加法与乘法运算。

这正是我们此前一直在处理的空间。
\end{example}

\subsection{多项式空间}

\begin{example}
$\mathcal{P}_n$ = 次数不超过 $n$ 的多项式全体:
\[
p(x) = a_0 + a_1 x + a_2 x^2 + \cdots + a_n x^n
\]

加法: $(p + q)(x) = p(x) + q(x)$

标量乘法: $(cp)(x) = c \cdot p(x)$

零向量: $p(x) = 0$

\textbf{基:} $\{1, x, x^2, \ldots, x^n\}$——空间维数为 $n + 1$
\end{example}

\subsection{函数空间}

\begin{example}
$C[a,b]$ = 区间 $[a,b]$ 上连续函数全体

该空间是\textbf{无穷维}的!
\end{example}

\subsection{矩阵空间}

\begin{example}
$M_{m \times n}$ = 全体 $m \times n$ 矩阵

加法: 按对应元相加

标量乘法: 所有元都乘以该标量

维数: $m \times n$
\end{example}

% ============================================
\section{抽象线性变换}
% ============================================

\begin{definition}[空间之间的线性变换]
函数 $T: V \to W$ 是\vocab{线性变换 (linear transformation)}, 如果:
\begin{enumerate}
    \item $T(\vu + \vv) = T(\vu) + T(\vv)$
    \item $T(c\vv) = c \cdot T(\vv)$
\end{enumerate}
\end{definition}

\begin{example}
\textbf{微分:} $D: \mathcal{P}_n \to \mathcal{P}_{n-1}$, $D(p) = p'$

这是线性的, 因为 $(f + g)' = f' + g'$ 且 $(cf)' = cf'$。
\end{example}

\begin{example}
\textbf{定积分:} $I: C[0,1] \to \R$, $I(f) = \int_0^1 f(x) dx$

这是线性的, 因为 $\int(f+g) = \int f + \int g$ 且 $\int cf = c\int f$。
\end{example}

% ============================================
\section{函数即无穷维向量}
% ============================================

\begin{intuition}
函数 $f: [0, 2\pi] \to \R$ 可视为一个无穷维向量:
\begin{itemize}
    \item 每个点 $x$ 是一个``分量''
    \item 函数值 $f(x)$ 就是该分量的取值
\end{itemize}

函数的内积:
\[
\langle f, g \rangle = \int_0^{2\pi} f(x) g(x) dx
\]
\end{intuition}

\begin{practical}
\textbf{傅里叶级数}

正弦与余弦函数构成周期函数的``基'':
\[
f(x) = \frac{a_0}{2} + \sum_{n=1}^{\infty} (a_n \cos nx + b_n \sin nx)
\]

系数 $a_n, b_n$ 就像该函数在这组基下的``坐标''!
\end{practical}

% ============================================
\section{抽象为何重要?}
% ============================================

\begin{summary}
\textbf{抽象的威力:}

一旦某个对象是向量空间, \textbf{线性代数的所有工具}都适用:
\begin{itemize}
    \item 线性无关 (linear independence) 与线性相关
    \item 基与维数
    \item 线性变换与矩阵
    \item 特征值与特征向量
    \item 像空间与零空间
\end{itemize}

线性代数中的一条定理 = 对所有这些空间都成立的定理!
\end{summary}

\begin{practical}
\textbf{量子力学}

量子态构成向量空间 (希尔伯特空间)。物理算符是线性变换。测量值 = 特征值!
\end{practical}

\begin{practical}
\textbf{信号处理}

音频信号是函数空间中的向量。傅里叶变换是基变换。滤波器是线性变换。
\end{practical}

\begin{practical}
\textbf{机器学习}

数据是特征空间中的向量。线性模型是线性变换。降维 = 寻找更好的基。
\end{practical}

% ============================================
\section{展望}
% ============================================

\begin{remark}
本课程只是线性代数的入门。更深入的主题包括:
\begin{itemize}
    \item \textbf{二次型}与二次曲线的分类
    \item \textbf{奇异值分解 (SVD)}——数据科学中的有力工具
    \item \textbf{数值线性代数}——高效算法
    \item \textbf{希尔伯特空间}——无穷维线性代数
    \item \textbf{表示论}——群与线性代数
\end{itemize}
\end{remark}

% ============================================
\section{练习}
% ============================================

\begin{exercise}
证明全体 $2 \times 2$ 对称矩阵构成一个向量空间。它的维数是多少?
\end{exercise}

\begin{exercise}
满足 $p(0) = 1$ 的多项式全体是否为向量空间? 为什么?
\end{exercise}

\begin{exercise}
证明微分 $D: \mathcal{P}_3 \to \mathcal{P}_2$ 是线性变换, 并在标准基下写出它的矩阵。
\end{exercise}

\begin{exercise}
证明在任一向量空间中: $0 \cdot \vv = \vzero$
\end{exercise}

\begin{exercise}[挑战]
考虑微分方程 $y'' + y = 0$ 的解空间。证明这是一个向量空间, 并求它的一组基。
\end{exercise}

\begin{problem}
对于线性变换 $T: \mathcal{P}_2 \to \mathcal{P}_2$, $T(p) = p + p'$:
\begin{enumerate}[label=(\alph*)]
    \item 在基 $\{1, x, x^2\}$ 下写出 $T$ 的矩阵
    \item 求特征值
    \item 求特征多项式 (特征向量)
\end{enumerate}
\end{problem}

% ============================================
\section{课程总结}
% ============================================

\begin{summary}
\textbf{线性代数的本质}

我们从向量与矩阵出发, 最终抵达抽象空间。核心思想:

\begin{enumerate}
    \item \textbf{向量:} 具有加法与标量乘法的对象
    \item \textbf{线性变换:} 保持直线不变的函数
    \item \textbf{矩阵:} 线性变换的数值表示
    \item \textbf{行列式:} 体积缩放因子
    \item \textbf{特征值:} 只被缩放的特殊方向
    \item \textbf{抽象:} 所有这些概念都超越箭头而成立
\end{enumerate}

线性代数是数学、物理学、工程学与计算机科学的共同语言。
\end{summary}


% === 附录: 术语表 ===
\appendix
\chapter{术语表}

\begin{center}
\begin{tabular}{l|l|l}
\toprule
\textbf{中文} & \textbf{英文} & \textbf{符号/记法} \\
\midrule
向量 & Vector & $\vv, \vw, \vu$ \\
矩阵 & Matrix & $\mA, \mB, \mC$ \\
线性变换 & Linear Transformation & $T: \Rtwo \to \Rtwo$ \\
行列式 & Determinant & $\det(\mA)$ \\
特征值 & Eigenvalue & $\lambda$ \\
特征向量 & Eigenvector & $\vv$ \\
基 & Basis & $\{\vi, \vj\}$ \\
张成 & Span & $\spn\{\vv_1, \vv_2, \ldots\}$ \\
零空间 & Null Space & $\Null(\mA)$ \\
列空间 & Column Space & $\Col(\mA)$ \\
行空间 & Row Space & $\Row(\mA)$ \\
秩 & Rank & $\rank(\mA)$ \\
线性无关 & Linear Independence & 无冗余的向量组 \\
线性相关 & Linear Dependence & 至少有一个冗余向量 \\
点积 & Dot Product & $\vv \cdot \vw$ \\
叉积 & Cross Product & $\vv \times \vw$ \\
单位矩阵 & Identity Matrix & $\mI$ \\
逆矩阵 & Inverse Matrix & $\mA\inv$ \\
转置 & Transpose & $\mA\trans$ \\
基变换 & Change of Basis & $\mP\inv \mA \mP$ \\
特征基 & Eigenbasis & 由特征向量组成的基 \\
对角矩阵 & Diagonal Matrix & 仅对角线非零 \\
标量 & Scalar & 一个数 (实数或复数) \\
线性组合 & Linear Combination & $c_1\vv_1 + c_2\vv_2 + \cdots$ \\
向量空间 & Vector Space & 对加法和标量乘法封闭的集合 \\
\bottomrule
\end{tabular}
\end{center}

\end{document}
